% =====================================================================
%  SYMPHONY V2 - Guida al layout della parte RF (UHF 435 MHz)
%  Riferita allo schematico rev. 1.45 del 04/11/2026
%  Compilare con pdflatex (due passaggi per indice e riferimenti)
% =====================================================================
\documentclass[a4paper,10pt]{article}

% ---------- lingua e font ----------
\usepackage[T1]{fontenc}
\usepackage[italian]{babel}
\usepackage[sups]{XCharter}
\usepackage[xcharter,vvarbb]{newtxmath}
\usepackage[scaled=0.88]{helvet}
\usepackage[scale=0.88,varqu]{inconsolata}
\usepackage[babel]{microtype}
\usepackage{csquotes}

% ---------- pagina ----------
\usepackage[a4paper,top=24mm,bottom=24mm,left=22mm,right=22mm,headheight=14pt]{geometry}
\usepackage{fancyhdr}
\usepackage{titlesec}
\usepackage{needspace}
\usepackage{etoolbox}
\pretocmd{\section}{\needspace{10\baselineskip}}{}{}
\pretocmd{\subsection}{\needspace{6\baselineskip}}{}{}
\usepackage{enumitem}
\usepackage{xcolor}
\usepackage{graphicx}
\usepackage{array,booktabs,tabularx,xltabular,multirow}
\usepackage{siunitx}
\usepackage{amsmath}
\usepackage{tikz}
\usetikzlibrary{arrows.meta,positioning,calc,patterns,decorations.pathreplacing,fit,backgrounds}
\usepackage[siunitx,RPvoltages]{circuitikz}
\usepackage{pgfplots}
\pgfplotsset{compat=1.18}
\usepackage[most]{tcolorbox}
\usepackage{pifont}
\usepackage{lastpage}
\usepackage{xurl}
\usepackage[hidelinks,bookmarksnumbered,pdfusetitle]{hyperref}
\urlstyle{tt}

\hypersetup{pdftitle={SYMPHONY V2 - Guida al layout della parte RF},
            pdfauthor={SYMPHONY V2 - documentazione hardware},
            pdfsubject={Layout RF UHF 435 MHz, schematico rev. 1.45}}

% ---------- colori ----------
\definecolor{sbblue}{HTML}{1F4E79}
\definecolor{sbgray}{HTML}{555555}
\definecolor{rxblue}{HTML}{1565C0}
\definecolor{txred}{HTML}{C62828}
\definecolor{rulegreen}{HTML}{2E7D32}
\definecolor{rulegreenbg}{HTML}{EEF6EE}
\definecolor{warnor}{HTML}{C75000}
\definecolor{warnorbg}{HTML}{FFF3E6}
\definecolor{whygray}{HTML}{5A6B7B}
\definecolor{whygraybg}{HTML}{F2F4F7}
\definecolor{verbg}{HTML}{FFF8D6}
\definecolor{verfr}{HTML}{9A7B00}

% ---------- siunitx ----------
\sisetup{output-decimal-marker={,},group-minimum-digits=5,
         per-mode=symbol,range-phrase={\,--\,},list-final-separator={ e }}

% ---------- titoli ----------
\titleformat{\section}{\sffamily\Large\bfseries\color{sbblue}}{\thesection}{0.8em}{}[\vspace{-0.6ex}{\color{sbblue}\rule{\linewidth}{0.6pt}}]
\titleformat{\subsection}{\sffamily\large\bfseries\color{sbblue}}{\thesubsection}{0.7em}{}
\titleformat{\subsubsection}{\sffamily\normalsize\bfseries\color{sbgray}}{\thesubsubsection}{0.6em}{}
\titlespacing*{\section}{0pt}{3.2ex plus 1ex minus .2ex}{1.6ex}
\titlespacing*{\subsection}{0pt}{2.6ex plus 1ex minus .2ex}{1.0ex}
\setcounter{tocdepth}{2}
\setcounter{secnumdepth}{3}

% ---------- intestazioni ----------
\pagestyle{fancy}
\fancyhf{}
\fancyhead[L]{\sffamily\footnotesize\color{sbgray}SYMPHONY V2 \textbullet{} Guida al layout RF}
\fancyhead[R]{\sffamily\footnotesize\color{sbgray}\nouppercase{\leftmark}}
\fancyfoot[L]{\sffamily\footnotesize\color{sbgray}Schematico rev.~1.45 (04/11/2026) \textbullet{} Guida rev.~A}
\fancyfoot[R]{\sffamily\footnotesize\color{sbgray}Pagina \thepage{} di \pageref*{LastPage}}
\renewcommand{\headrulewidth}{0.4pt}
\renewcommand{\footrulewidth}{0.4pt}
\fancypagestyle{plain}{\fancyhf{}\renewcommand{\headrulewidth}{0pt}\renewcommand{\footrulewidth}{0pt}}

% ---------- liste ----------
\setlist{itemsep=2pt,topsep=4pt,parsep=0pt}
\setlist[itemize]{leftmargin=1.4em}
\setlist[enumerate]{leftmargin=1.8em}
\newlist{checklist}{itemize}{2}
\setlist[checklist]{label=\Large$\square$,leftmargin=2.0em,itemsep=3pt}

% ---------- comandi ----------
\newcommand{\net}[1]{\texttt{\detokenize{#1}}}
\newcommand{\cmp}[1]{\textsf{#1}}
\newcommand{\dv}{\textcolor{verfr}{\textbf{[da verificare]}}}
\newcommand{\fonte}[1]{{\footnotesize\color{sbgray}[#1]}}
\newcolumntype{L}[1]{>{\raggedright\arraybackslash}p{#1}}
\newcolumntype{Y}{>{\raggedright\arraybackslash}X}

% ---------- riquadri ----------
\newcounter{regola}
\newtcolorbox{regola}[1][]{enhanced,breakable,colback=rulegreenbg,colframe=rulegreen,
  boxrule=0.8pt,arc=1.5pt,left=6pt,right=6pt,top=4pt,bottom=4pt,
  fonttitle=\sffamily\bfseries\small,coltitle=white,colbacktitle=rulegreen,
  title={\stepcounter{regola}Regola R\theregola\ifx\relax#1\relax\else{} \textendash{} #1\fi}}
\newtcolorbox{attenzione}[1][]{enhanced,breakable,colback=warnorbg,colframe=warnor,
  boxrule=0.8pt,arc=1.5pt,left=6pt,right=6pt,top=4pt,bottom=4pt,
  fonttitle=\sffamily\bfseries\small,coltitle=white,colbacktitle=warnor,
  title={Attenzione\ifx\relax#1\relax\else{} \textendash{} #1\fi}}
\newtcolorbox{perche}[1][]{enhanced,breakable,colback=whygraybg,colframe=whygray,
  boxrule=0.6pt,arc=1.5pt,left=6pt,right=6pt,top=4pt,bottom=4pt,
  fonttitle=\sffamily\bfseries\small,coltitle=white,colbacktitle=whygray,
  title={Perché\ifx\relax#1\relax\else{} \textendash{} #1\fi}}
\newtcolorbox{verifica}[1][]{enhanced,breakable,colback=verbg,colframe=verfr,
  boxrule=0.6pt,arc=1.5pt,left=6pt,right=6pt,top=4pt,bottom=4pt,
  fonttitle=\sffamily\bfseries\small,coltitle=white,colbacktitle=verfr,
  title={Da verificare\ifx\relax#1\relax\else{} \textendash{} #1\fi}}

% ---------- stili TikZ ----------
\tikzset{
  blk/.style={draw=black!70,rounded corners=1.2pt,line width=0.6pt,align=center,
              font=\sffamily\scriptsize,inner sep=2pt},
  rxb/.style={blk,fill=rxblue!10},
  txb/.style={blk,fill=txred!10},
  pab/.style={blk,fill=txred!22},
  sib/.style={blk,fill=green!10},
  ckb/.style={blk,fill=yellow!25},
  neu/.style={blk,fill=black!6},
  rxl/.style={draw=rxblue,line width=1.4pt,-{Stealth[length=2.6mm]}},
  txl/.style={draw=txred,line width=1.4pt,-{Stealth[length=2.6mm]}},
  fence/.style={draw=black!65,line width=0.9pt,dash pattern=on 0.8pt off 1.6pt,line cap=round},
  dig/.style={draw=violet,line width=0.9pt,dashed,-{Stealth[length=2mm]}},
  dim/.style={draw=black!60,line width=0.4pt,{Latex[length=1.6mm]}-{Latex[length=1.6mm]}},
  lab/.style={font=\sffamily\tiny,text=black!75,align=center},
}

\begin{document}

% =====================================================================
%  FRONTESPIZIO
% =====================================================================
\pagenumbering{alph}
\begin{titlepage}
\thispagestyle{empty}
\begin{tikzpicture}[remember picture,overlay]
  \fill[sbblue] (current page.north west) rectangle ([yshift=-62mm]current page.north east);
  \node[anchor=north west,text=white,font=\sffamily\bfseries\fontsize{30}{34}\selectfont,align=left]
     at ([xshift=22mm,yshift=-18mm]current page.north west) {SYMPHONY V2};
  \node[anchor=north west,text=white!85!sbblue,font=\sffamily\Large,align=left]
     at ([xshift=22mm,yshift=-34mm]current page.north west) {Ricetrasmettitore UHF \SI{435}{\mega\hertz} per CubeSat};
  \fill[txred] ([yshift=-62mm]current page.north west) rectangle ([yshift=-64mm]current page.north east);
\end{tikzpicture}

\vspace*{62mm}
{\sffamily\bfseries\fontsize{26}{31}\selectfont\color{sbblue} Guida al layout\\[2pt] della parte RF\par}
\vspace{6mm}
{\sffamily\large\color{sbgray} Piazzamento, sbroglio, masse, stackup, netclass e controlli\\
per la catena antenna -- switch -- LNA -- Si4464 -- PA -- filtro\par}

\vspace{14mm}
\begin{tcolorbox}[colback=black!3,colframe=black!35,boxrule=0.5pt,arc=1pt,width=0.86\linewidth,
                  left=8pt,right=8pt,top=6pt,bottom=6pt]
\sffamily\small
\begin{tabular}{@{}L{44mm}L{86mm}@{}}
Documento & Guida al layout RF, rev.~A\\
Data & 04/11/2026\\
Riferimento di progetto & Schematico SYMPHONY V2 \textbf{rev.~1.45 del 04/11/2026}\\
                        & (netlist esportata con \texttt{kicad-cli} 10.0.8 il 04/11/2026)\\
Note di layout riprese & \S24 del log (rev.~1.12) e nota L-01 (rev.~1.45)\\
Stato del PCB & non ancora sbrogliato; stackup non ancora deciso\\
Destinatario & chi disegna il PCB (parte RF)\\
\end{tabular}
\end{tcolorbox}

\vfill
{\small\color{sbgray}
Questa guida tratta solo la parte RF del foglio UHF e i punti in cui il resto della scheda la tocca
(alimentazione del PA, masse, separazione dal digitale, TCXO). Ogni valore numerico ha una fonte
(datasheet o application note, con pagina) o un calcolo mostrato; quando la fonte non era verificabile
il valore è marcato \dv. Lo schematico e il file di progetto KiCad non sono stati modificati.\par}
\end{titlepage}

% =====================================================================
\pagenumbering{roman}
\tableofcontents
\clearpage
\pagenumbering{arabic}

% =====================================================================
%  SINTESI
% =====================================================================
\section*{Sintesi: le 10 regole più importanti}
\addcontentsline{toc}{section}{Sintesi: le 10 regole più importanti}
\markboth{Sintesi}{Sintesi}

{\small
\renewcommand{\arraystretch}{1.25}
\begin{tabularx}{\linewidth}{@{}>{\sffamily\bfseries\color{sbblue}}c Y L{30mm}@{}}
\toprule
\textnormal{\#} & \textbf{Regola} & \textbf{Dove}\\
\midrule
1 & \textbf{Piano di massa In1.Cu continuo} sotto tutta la zona RF: nessuna pista, nessun taglio. Sotto la zona RF anche In2.Cu e B.Cu solo GND, tranne un unico corridoio digitale (regola 10). & \S\ref{sec:stackup}, \S\ref{sec:generali}\\
2 & \textbf{Linee \SI{50}{\ohm} coplanari con massa (CPWG) su F.Cu}: pista \SI{0.20}{mm}, distanza dal GND \SI{0.20}{mm} sullo stackup JLC04161H-3313. Corte, dritte, senza via e senza stub. Confermare con il calcolatore JLCPCB. & \S\ref{sec:stackup}\\
3 & \textbf{Catena in linea, ad anello}: J201 sul bordo, C206, IC202; TX sopra (esce da RF2), RX sotto (esce da RF1). Per l'ordine dei pin il \textbf{Si4464 va dentro l'anello con il lato RF verso l'esterno}, altrimenti TX e RX si incrociano. & \S\ref{sec:floorplan}\\
4 & \textbf{LNA vicino allo switch}: da J201 all'ingresso di Q201 al massimo circa \SI{15}{mm} di pista, filtro C210--L201--C209 compreso. & \S\ref{sec:rx}\\
5 & \textbf{Rete d'uscita del PA = layout EVB Qorvo} (nota L-01): stesso ordine, stesse posizioni e lunghezze; C260, C261, C262 con via di massa propria sul pad (almeno 2 per C260 e C261); \SI{1}{mm} libero intorno a C260 e C262 per la taratura. & \S\ref{sec:pa}\\
6 & \textbf{PA: massa e calore}. Pad esposto con matrice di via, altre via nel GND intorno, rame su B.Cu verso la struttura. Pin~4 (GND2) a massa con pista di circa \SI{1}{mm} (40 mil), pin~2 diretto. Ramo 3V3$\rightarrow$L242: almeno \SI{0.62}{mm} (IPC-2221), consigliato \SI{1}{mm} o un pour, dedicato dal buck. & \S\ref{sec:pa}\\
7 & \textbf{FL201 tra PA e switch}, ingresso e uscita in direzioni opposte, via su tutti i 6 pad GND, fila di via trasversale sotto il filtro: l'uscita non deve \enquote{vedere} l'ingresso. & \S\ref{sec:fl201}\\
8 & \textbf{Rete AN643 compatta e simmetrica sui pin 2/3 del Si4464}: \net{RF_RX_P} e \net{RF_RX_N} sotto \SI{3}{mm} e uguali entro \SI{0.5}{mm}; L206 subito tra i pad RXP e RXN, C220 al pad RXN con via accanto. Footprint standard, non HandSolder. & \S\ref{sec:rx}\\
9 & \textbf{Via fence e cucitura}: passo $\le$\,\SI{2}{mm} lungo le linee RF e sul bordo (\SI{1.5}{mm} intorno a PA e FL201). Fila di via tra zona TX e zona RX. PA ad almeno \SI{15}{mm} da Y201 e dalla catena RX. & \S\ref{sec:generali}\\
10 & \textbf{TCXO vicino a XIN} (C204, R201 in linea, $\le$\,\SI{5}{mm} in tutto \dv), lontano da PA e L242. \textbf{Digitale mai sotto o vicino all'RF} ($\ge$\,\SI{2}{mm} su F.Cu); i segnali che entrano nell'anello lo attraversano una sola volta, a \ang{90}, su B.Cu/In2 sotto In1 GND, nel tratto TX a basso livello. & \S\ref{sec:tcxo}, \S\ref{sec:floorplan}\\
\bottomrule
\end{tabularx}}

\medskip
\begin{attenzione}[prima di iniziare]
Nel file \texttt{SYMPHONY.kicad\_pro} attuale \textbf{non ci sono le netclass RF}: c'è solo la classe
\texttt{Default}, nessuna assegnazione e nessun pattern, anche se il log (rev.~1.12 e 1.13) le dà come
inserite. Vanno ricreate prima di sbrogliare (\S\ref{sec:netclass}). Il file \texttt{SYMPHONY.kicad\_pcb}
esistente ha inoltre uno stackup a 6 strati con \texttt{In1.Cu} di segnale sotto F.Cu, incompatibile con
le regole di questa guida (\S\ref{sec:stackup}).
\end{attenzione}

\clearpage
% =====================================================================
\section{Scopo, fonti e convenzioni}
\label{sec:intro}

\subsection{Cosa copre questa guida}
La guida dice come piazzare e sbrogliare tutta la catena RF del foglio UHF: connettore d'antenna J201,
switch HMC349A (IC202), catena RX con LNA TQP3M9036 (Q201), filtri e rete AN643, ricetrasmettitore
Si4464 (IC204) con TCXO (Y201), rete TX AN648, attenuatore, PA QPA9510 (IC2) con la sua rete d'uscita,
filtro passa-basso LFCG-490+ (FL201). Il resto della scheda resta al progettista; qui se ne parla solo
dove tocca la parte RF.

Riprende e integra la guida al layout UHF del log (\S24, rev.~1.12) e la nota di layout L-01
(rev.~1.45). Dove questa guida aggiunge o precisa qualcosa rispetto a quelle note lo dice in modo
esplicito; le differenze da decidere sono raccolte nell'Appendice~\ref{app:verifiche}.

\subsection{Fonti}
\begin{itemize}
\item Netlist KiCad esportata il 04/11/2026 da \texttt{SYMPHONY.kicad\_sch} (rev.~1.45): componenti,
      valori, footprint e collegamenti di questa guida vengono da lì (file \texttt{net.xml} nella cartella
      di questo documento).
\item Log hardware \texttt{report/entries}: 001 (picchi del PA, alimentazione dedicata), 007 (rete RX
      AN643), 009 e 023 (L242), 010 (\S24, netclass e guida UHF), 043 (nota L-01).
\item Revisione indipendente del 04/11/2026 (\texttt{SYMPHONY\_V2\_review\_finale\_2026-11-04.pdf}), punti P05, P08, P10.
\item Datasheet: Si4464 rev.~1.4; QPA9510 del 23/06/2025 e \emph{QPA9510 410--470\,MHz EVB Test
      Summary} del 30/06/2025; HMC349AMS8G rev.~F; TQP3M9036 del 27/11/2018; Mini-Circuits LFCG-490+;
      ECS-TXO-20CSMV4; DreamLNK DL-MMCX-KEF-180.
\item JLCPCB: tabella degli stackup a 4 strati (\nolinkurl{jlcpcb.com/impedance}) e capacità di
      produzione (\nolinkurl{jlcpcb.com/capabilities/pcb-capabilities}), consultate il 04/11/2026.
\item Silicon Labs AN629 (layout Si446x), AN643 (match RX) e AN648 (match TX): il sito Silicon Labs ha
      rifiutato l'accesso ai PDF durante la stesura (HTTP 403). I valori di AN643 e AN648 sono quelli già
      riportati nel log (entries 001, 007, 012); le regole attribuite ad AN629 sono marcate \dv.
\end{itemize}

\subsection{Convenzioni}
\begin{regola}[esempio]
Un riquadro verde è una regola da seguire. Le regole sono numerate R1, R2, \ldots{} in tutto il documento.
\end{regola}
\addtocounter{regola}{-1}
\begin{attenzione}
Un riquadro arancio segnala un errore facile da fare o un punto aperto del progetto.
\end{attenzione}
\begin{perche}
Un riquadro grigio spiega il motivo di una regola, con il calcolo quando serve.
\end{perche}
Le reti sono scritte in \net{carattere_macchina} con il nome della netlist senza il prefisso
\net{/UHF/}. Le quote sono in millimetri. \enquote{GND} è la massa unica della scheda.

% =====================================================================
\section{La catena RF secondo la netlist}
\label{sec:catena}

\subsection{Percorso del segnale}
\begin{center}
\begin{tikzpicture}[node distance=3.5mm and 3.2mm, every node/.style={font=\sffamily\scriptsize}]
\node[neu] (j) {J201\\MMCX};
\node[neu,right=of j] (c206) {C206\\100p};
\node[neu,right=of c206,minimum height=12mm] (sw) {IC202\\HMC349A};
% RX
\node[rxb,below right=6mm and 3mm of sw] (c207) {C207\\100p};
\node[rxb,right=of c207] (lpf) {C210\,/\,L201\,/\,C209\\LPF $\pi$ 6,8p--22n--6,8p};
\node[rxb,right=of lpf] (c212) {C212\\100p};
\node[rxb,right=of c212] (lna) {Q201\\TQP3M9036\\(L202 bias)};
\node[rxb,below=6mm of lna] (f2) {C213 + C217\,/\,L203\,/\\C216\,/\,L204\,/\,C218};
\node[rxb,left=of f2] (an643) {C219, L205,\\L206, C220\\AN643};
\node[sib,left=of an643,minimum height=12mm] (si) {IC204\\Si4464};
\node[ckb,below=4mm of si] (y) {Y201 TCXO 32\,MHz\\C204, R201 $\to$ XIN};
% TX
\node[txb,above right=6mm and 3mm of sw] (c208) {C208\\100p};
\node[txb,right=of c208] (fl) {FL201\\LFCG-490+};
\node[txb,right=of fl] (c230) {C230\\56p};
\node[pab,right=of c230] (om) {C261/C262, L243,\\C260, L242};
\node[pab,right=of om] (pa) {IC2\\QPA9510};
\node[txb,above=6mm of pa] (att) {C255 + $\pi$\\R245--R248};
\node[txb,left=of att] (an648) {C701, C704, L702,\\C702, L703\\AN648};
% wires
\draw[line width=1pt] (j) -- (c206) -- (sw);
\draw[rxl] (sw.south east) |- (c207.west);
\draw[rxl] (c207) -- (lpf); \draw[rxl] (lpf) -- (c212); \draw[rxl] (c212) -- (lna);
\draw[rxl] (lna) -- (f2); \draw[rxl] (f2) -- (an643); \draw[rxl] (an643) -- (si);
\draw[draw=olive,line width=0.9pt,-{Stealth[length=2mm]}] (y) -- (si);
\draw[txl] (si.north) |- (an648.west);
\draw[txl] (an648) -- (att); \draw[txl] (att) -- (pa); \draw[txl] (pa) -- (om);
\draw[txl] (om) -- (c230); \draw[txl] (c230) -- (fl); \draw[txl] (fl) -- (c208);
\draw[txl] (c208.west) -| (sw.north east);
\node[anchor=west,text=rxblue,font=\sffamily\scriptsize\bfseries] at ($(c207.south west)+(0,-3mm)$) {RX: RF1 (pin 5)};
\node[anchor=west,text=txred,font=\sffamily\scriptsize\bfseries] at ($(c208.north west)+(0,3mm)$) {TX: RF2 (pin 8)};
\end{tikzpicture}
\end{center}
\noindent{\footnotesize Schema a blocchi della netlist rev.~1.45 (non è un piazzamento). In blu la catena RX, in
rosso la TX.}

\subsection{Reti RF con nome}
{\footnotesize
\renewcommand{\arraystretch}{1.12}
\begin{xltabular}{\linewidth}{@{}L{27mm}L{62mm}Y@{}}
\toprule
\textbf{Rete} & \textbf{Nodi (netlist rev.~1.45)} & \textbf{Ruolo per il layout}\\
\midrule\endhead
\midrule\multicolumn{3}{r}{\footnotesize\itshape continua\ldots}\\\endfoot
\bottomrule\endlastfoot
\net{RF_ANT} & J201.1, C206.2 & linea d'antenna \SI{50}{\ohm}, TX fino a circa \SI{30}{dBm}\\
\net{RF_SW_COM} & C206.1, IC202.3 (RFC) & porta comune dello switch\\
\net{RF_SW_RX} & IC202.5 (RF1), C207.2 & uscita RX dello switch\\
\net{RF_SW_TX} & IC202.8 (RF2), C208.2 & ingresso TX dello switch\\
\net{RF_RX_F1}, \net{RF_RX_F2} & C207, C210, L201 / L201, C209, C212 & filtro passa-basso $\pi$ prima del LNA\\
\net{RF_LNA_IN} & C212.1, Q201.2 & ingresso LNA (punto più sensibile)\\
\net{RF_LNA_OUT} & Q201.7, C213.2, L202.1 & uscita LNA + bias (L202 da 3V3)\\
\net{RF_RX_F3}, \net{RF_RX_F4} & C213, C217, L203 / L203, C216, L204 & filtro dopo il LNA\\
\net{RF_RX_MATCH} & L204, C218, C219, L205 & uscita filtro, ingresso rete AN643\\
\net{RF_RX_P} & C219.1, IC204.2 (RXP), L206.1 & ingresso differenziale, alta impedenza\\
\net{RF_RX_N} & L205.1, L206.2, C220.2, IC204.3 (RXN) & ingresso differenziale, alta impedenza\\
\net{RF_SI_TX} & IC204.4 (TX), L703.2, C702.2 & pin TX in classe E (picchi di tensione)\\
\net{RF_TX_M1}, \net{RF_TX_M2} & C702, L702 / L702, C704, C701 & match AN648\\
\net{RF_ATT_IN}, \net{RF_ATT_OUT} & C701, R245, R246 / R246, R247, R248, C255 & attenuatore $\pi$ (\SI{10.8}{dB})\\
\net{RF_PA_IN} & C255.2, IC2.3 & ingresso PA\\
\net{RF_PA_OUT} & IC2.9--12, C260, L243, L242 & nodo a circa \SI{5}{\ohm}, \SI{1.7}{A} DC\\
\net{RF_PA_MATCH} & L243, C261, C262, C230 & nodo intermedio del match d'uscita\\
\net{RF_LPF_IN}, \net{RF_LPF_OUT} & C230, FL201.8 / FL201.4, C208 & ingresso e uscita del filtro\\
\net{RF_TCXO_OUT}, \net{RF_TCXO_AC}, \net{RF_XIN} & Y201.3, C204 / C204, R201 / R201, IC204.17, TP701 & clock \SI{32}{MHz} al Si4464\\
\net{PA_VCC1}, \net{PA_VCC1_F} & IC2.1, L240 / L240, R240, C244 & alimentazione e match interstadio\\
\net{PA_VCC2}, \net{PA_VCC2_F} & IC2.5--7, L241 / L241, R241, C248 & alimentazione e match interstadio\\
\end{xltabular}}

Il resto delle connessioni che servono al layout: IC2 pin~14 (VCC) su 3V3 con C251; APC1/APC2
(pin 15, 16) sulla rete \net{Net-(IC2-APC1)} con R242, R243, C252; pin~8 (2F0) e pin~13 (NC)
scollegati; pin~2, 4 e pad esposto 17 a GND. IC202: VDD (pin~1) su 5V0 con C211, CTRL (pin~2) da
\net{EN_TX_ANT_SW}, EN (pin~4) a GND. Q201: SHUT\_DOWN (pin~6) da \net{EN_RX_LNA} tramite R202 (R203 DNP
verso GND), pin NC 1, 3, 4, 5, 8 a GND. IC204: VDD pin~6 e 8 su 3V3 (C201), XOUT (pin~16) scollegato.

% =====================================================================
\section{Stackup e linee a 50 ohm}
\label{sec:stackup}

\subsection{Situazione}
Lo stackup non è ancora deciso. Il log (\S24.1) indica un 4 strati da \SI{1.6}{mm} con prepreg 3313 e
In1.Cu a GND; il file \texttt{SYMPHONY.kicad\_pcb} attuale contiene invece uno stackup a 6 strati
(F.Cu, In1 \enquote{SIG}, In2 \enquote{GND}, In3 \enquote{SIG}, In4 \enquote{POWER}, B.Cu), con un piano
di segnale subito sotto F.Cu. Questa guida assume il 4 strati JLCPCB qui sotto.

\begin{attenzione}[stackup del file PCB]
Con lo stackup a 6 strati presente oggi nel file PCB, le linee RF su F.Cu avrebbero come riferimento
In1 \enquote{SIG}: non va bene. Se si resta a 6 strati, In1 deve diventare GND pieno sotto la zona RF e
le larghezze vanno ricalcolate.
\end{attenzione}

\subsection{Stackup proposto: JLCPCB JLC04161H-3313}

\begin{center}
\begin{minipage}{0.46\linewidth}
\centering
\begin{tikzpicture}[x=1mm,y=1mm]
\def\w{52}
% layers from top (y decreasing)
\fill[green!45!black!60] (0,40) rectangle (\w,41) node[midway,text=white,font=\sffamily\tiny]{solder mask};
\fill[orange!80!black] (0,38.6) rectangle (\w,40);
\node[anchor=west,font=\sffamily\tiny] at (\w+1,39.3) {F.Cu \SI{35}{\micro\metre}: RF, componenti};
\fill[yellow!25] (0,33.6) rectangle (\w,38.6);
\node[font=\sffamily\tiny] at (\w/2,36.1) {prepreg 3313, \SI{0.0994}{mm}, $\varepsilon_r$ 4,1};
\fill[orange!80!black] (0,32.8) rectangle (\w,33.6);
\node[anchor=west,font=\sffamily\tiny] at (\w+1,33.2) {In1.Cu \SI{15}{\micro\metre}: \textbf{GND pieno}};
\fill[yellow!45] (0,14.2) rectangle (\w,32.8);
\node[font=\sffamily\tiny,align=center] at (\w/2,23.5) {core \SI{1.265}{mm}, $\varepsilon_r$ 4,6\\(non in scala)};
\fill[orange!80!black] (0,13.4) rectangle (\w,14.2);
\node[anchor=west,font=\sffamily\tiny] at (\w+1,13.8) {In2.Cu \SI{15}{\micro\metre}: alimentazioni};
\fill[yellow!25] (0,8.4) rectangle (\w,13.4);
\node[font=\sffamily\tiny] at (\w/2,10.9) {prepreg 3313, \SI{0.0994}{mm}};
\fill[orange!80!black] (0,7) rectangle (\w,8.4);
\node[anchor=west,font=\sffamily\tiny] at (\w+1,7.7) {B.Cu \SI{35}{\micro\metre}: segnali, GND};
\fill[green!45!black!60] (0,6) rectangle (\w,7);
\end{tikzpicture}
\end{minipage}\hfill
\begin{minipage}{0.48\linewidth}
\small
\begin{tabular}{@{}ll@{}}
\toprule
\textbf{Strato} & \textbf{Spessore / materiale}\\
\midrule
F.Cu & \SI{0.035}{mm} (1 oz)\\
prepreg & 3313, \SI{0.0994}{mm}, $\varepsilon_r=4{,}1$\\
In1.Cu & \SI{0.0152}{mm} (0,5 oz)\\
core & \SI{1.265}{mm}, $\varepsilon_r=4{,}6$\\
In2.Cu & \SI{0.0152}{mm} (0,5 oz)\\
prepreg & 3313, \SI{0.0994}{mm}, $\varepsilon_r=4{,}1$\\
B.Cu & \SI{0.035}{mm} (1 oz)\\
\midrule
totale & circa \SI{1.56}{mm} + solder mask\\
\bottomrule
\end{tabular}

\smallskip
{\footnotesize\color{sbgray}Fonte: JLCPCB, tabella \enquote{4 layers} di \texttt{jlcpcb.com/impedance}
(consultata il 04/11/2026).}
\end{minipage}
\end{center}

Uso degli strati nella zona RF: F.Cu per linee RF e componenti; In1.Cu \textbf{solo GND, senza
interruzioni}; In2.Cu alimentazioni (3V3, 5V0, pour del PA), ma sotto le linee RF solo GND; B.Cu segnali
lenti e GND, sotto la zona RF solo GND cucito.

\begin{attenzione}[spessori nel log]
Il \S24.1 del log scrive \enquote{F.Cu -- prepreg 0.1 -- In1 -- core 0.535 -- In2 -- prepreg 0.1 -- B.Cu,
\SI{1.6}{mm}}: la somma fa circa \SI{0.88}{mm}, non \SI{1.6}{mm}. Il core del JLC04161H-3313 è
\SI{1.265}{mm}. Per le linee su F.Cu conta solo il prepreg superiore (\SI{0.0994}{mm}), quindi le
larghezze del \S24.1 restano valide.
\end{attenzione}

\subsection{Formule usate}
\textbf{Microstriscia} (Hammerstad--Jensen, con correzione per lo spessore $t$ del rame). Con $u=W/h$:
\begin{gather*}
Z_{01}(u)=\frac{\eta_0}{2\pi}\ln\!\left[\frac{f(u)}{u}+\sqrt{1+\frac{4}{u^2}}\right],\qquad
f(u)=6+(2\pi-6)\,e^{-(30{,}666/u)^{0{,}7528}}\\[2pt]
\varepsilon_\mathrm{eff}(u)=\frac{\varepsilon_r+1}{2}+\frac{\varepsilon_r-1}{2}\left(1+\frac{10}{u}\right)^{-ab},
\qquad Z_0=\frac{Z_{01}(u_r)}{\sqrt{\varepsilon_\mathrm{eff}(u_r)}}
\end{gather*}
dove $a$, $b$ sono i coefficienti di Hammerstad--Jensen e $u_r=u+\Delta u_r$ è la larghezza corretta per
lo spessore. $\eta_0=\SI{376.73}{\ohm}$.

\textbf{Coplanare con massa (CPWG)}, metodo della trasformazione conforme, spessore nullo. Con pista $W$,
distanza $s$ dal GND di F.Cu e dielettrico $h$ verso In1:
\begin{gather*}
k=\frac{W}{W+2s},\qquad k_3=\frac{\tanh(\pi W/4h)}{\tanh\!\big(\pi (W+2s)/4h\big)},\qquad k'=\sqrt{1-k^2}\\[2pt]
\varepsilon_\mathrm{eff}=\frac{1+\varepsilon_r\,\dfrac{K(k')}{K(k)}\dfrac{K(k_3)}{K(k_3')}}{1+\dfrac{K(k')}{K(k)}\dfrac{K(k_3)}{K(k_3')}},\qquad
Z_0=\frac{60\pi}{\sqrt{\varepsilon_\mathrm{eff}}}\;\frac{1}{\dfrac{K(k)}{K(k')}+\dfrac{K(k_3)}{K(k_3')}}
\end{gather*}
con $K(k)/K(k')$ calcolato con l'approssimazione di Hilberg (errore $<10^{-5}$).

Il calcolo non tiene conto della solder mask, che abbassa $Z_0$ di 1--3\,\si{\ohm} su piste così strette
\dv; la CPWG è a spessore nullo, e lo spessore reale (\SI{35}{\micro\metre}) la abbassa ancora di poco.
Per questo i risultati vanno confermati con il calcolatore JLCPCB.

\subsection{Risultati}
{\small
\begin{center}
\begin{tabular}{@{}lccccc@{}}
\toprule
\textbf{Stackup JLC} & $h$ [mm] & $\varepsilon_r$ & \textbf{Microstriscia} & \textbf{CPWG $s$\,=\,0,20} & \textbf{CPWG $s$\,=\,0,15}\\
& & & $W_{50}$ [mm] ($\varepsilon_\mathrm{eff}$) & $W_{50}$ [mm] ($\varepsilon_\mathrm{eff}$) & $W_{50}$ [mm]\\
\midrule
\textbf{JLC04161H-3313} & 0,0994 & 4,10 & \textbf{0,175} (2,93) & \textbf{0,206} (3,16) & 0,200\\
JLC04161H-2116 & 0,1164 & 4,16 & 0,206 (2,98) & 0,235 (3,16) & 0,227\\
JLC04161H-1080 & 0,0764 & 3,91 & 0,135 (2,78) & 0,168 (3,09) & 0,163\\
JLC04161H-7628 & 0,2104 & 4,40 & 0,372 (3,19) & 0,376 (3,17) & 0,355\\
\bottomrule
\end{tabular}
\end{center}}
\noindent{\footnotesize Rame esterno \SI{35}{\micro\metre}. Verifiche puntuali sul 3313: CPWG \SI{0.20}{}/\SI{0.20}{mm}
$=\SI{50.9}{\ohm}$ con $\varepsilon_r=4{,}1$ e \SI{48.8}{\ohm} con $\varepsilon_r=4{,}5$ (stesso risultato del
\S24.1); microstriscia \SI{0.20}{mm} $=\SI{46.5}{\ohm}$; microstriscia \SI{0.5}{mm} $=\SI{25.6}{\ohm}$;
\SI{1.0}{mm} $=\SI{14.9}{\ohm}$.}

\begin{verifica}[impedenze]
Verificare tutte le larghezze con il calcolatore d'impedenza JLCPCB per lo stackup scelto, con solder
mask, prima di sbrogliare. Se la larghezza cambia, aggiornare la netclass \net{RF_50R}.
\end{verifica}

\subsection{Quale linea e perché}

\begin{regola}[tipo di linea]
Usare la \textbf{CPWG}: pista \SI{0.20}{mm}, distanza dal GND di F.Cu \SI{0.20}{mm}, In1.Cu GND pieno
sotto, via di cucitura lungo i due bordi del GND (\S\ref{sec:generali}). È la scelta del \S24.1.
\end{regola}

\begin{perche}[CPWG invece di microstriscia]
\begin{itemize}
\item Con un dielettrico di \SI{0.1}{mm} la CPWG da \SI{50}{\ohm} è quasi una microstriscia (il GND di F.Cu
      a \SI{0.2}{mm} sposta poco l'impedenza: \SI{0.175}{mm} contro \SI{0.206}{mm}), ma porta il GND di
      superficie fino al bordo della linea. Così i condensatori in parallelo trovano massa e via a pochi
      decimi di millimetro e si può mettere la fila di via a ridosso della linea.
\item È la stessa linea usata nella EVB dell'HMC349A (CPWG, \SI{0.41}{mm} di pista e \SI{0.33}{mm} di
      distanza su \SI{0.25}{mm} di RO4350, \fonte{HMC349AMS8G DS rev.~F p.~9}) e della TQP3M9036
      (\SI{0.51}{mm}/\SI{0.81}{mm} su \SI{0.25}{mm} RO4350B, \fonte{TQP3M9036 DS p.~10}).
\item Una pista più stretta del pad (\SI{0.2}{mm} contro \SIrange{0.5}{0.9}{mm} dei pad 0402--0805) crea un
      gradino. A \SI{435}{MHz} su linee da \SI{50}{\ohm} è trascurabile: il pad di un 0603
      (\SI{1.0}{}$\times$\SI{0.6}{mm}) sopra \SI{0.1}{mm} di prepreg vale circa
      $C=\varepsilon_0\varepsilon_r A/h\approx\SI{0.22}{pF}$, cioè $|X|\approx\SI{1.7}{k\ohm}$. Diventa invece
      importante nei punti ad alta impedenza (ingresso del Si4464) o a bassissima impedenza (uscita del PA),
      trattati a parte.
\end{itemize}
\end{perche}

\begin{perche}[perché il prepreg sottile (3313) va bene]
\begin{itemize}
\item La via di massa di un condensatore in parallelo arriva subito al piano: l'induttanza del tratto
      F.Cu--In1 vale circa $L\approx 0{,}2\,h_v\left[\ln(4h_v/d)+1\right]\,\si{nH}$ ($h_v$, $d$ in mm),
      con $h_v\approx\SI{0.13}{mm}$ e $d=\SI{0.3}{mm}$: \SI{0.04}{nH} ($|X|\approx\SI{0.11}{\ohm}$ a
      \SI{435}{MHz}). Con il 7628 ($h_v\approx\SI{0.25}{mm}$) diventa \SI{0.11}{nH}.
\item Linee strette e campo confinato: meno accoppiamento tra piste vicine (\S\ref{sec:generali}).
\end{itemize}
Lo svantaggio è nella rete d'uscita del PA: la EVB Qorvo usa piste di circa \SI{0.5}{mm}, che sul 3313
valgono circa \SI{26}{\ohm} invece di \SI{50}{\ohm} (\S\ref{sec:pa}). Con il JLC04161H-7628 la stessa pista
da \SI{0.5}{mm} vale \SI{42}{\ohm}, più vicino alla EVB. È un'alternativa da valutare, non la scelta di
questa guida.
\end{perche}

\subsection{Perdite e lunghezza d'onda a 435 MHz}
\begin{itemize}
\item $\lambda_0=c/f=\SI{689.7}{mm}$; in linea CPWG ($\varepsilon_\mathrm{eff}\approx3{,}16$)
      $\lambda_g=\lambda_0/\sqrt{\varepsilon_\mathrm{eff}}\approx\SI{388}{mm}$. Una pista di \SI{10}{mm} è
      $\lambda_g/39$: corta, ma non trascurabile nei nodi a bassa impedenza.
\item Perdita stimata della linea da \SI{0.2}{mm} su FR4 ($\tan\delta=0{,}02$ \dv):
      dielettrico
      \[\alpha_d\approx27{,}3\,\frac{\varepsilon_r}{\sqrt{\varepsilon_\mathrm{eff}}}\,\frac{\varepsilon_\mathrm{eff}-1}{\varepsilon_r-1}\,\frac{\tan\delta}{\lambda_0}\approx\SI{0.0012}{dB/mm};\]
      conduttore $\alpha_c\approx8{,}686\,R_s/(Z_0W)$ con $R_s=\sqrt{\pi f\mu_0\rho}=\SI{5.4}{m\ohm}$:
      $\approx\SI{0.0047}{dB/mm}$. In tutto circa \SI{0.006}{dB/mm}, cioè \SI{0.06}{dB} ogni \SI{10}{mm}.
\end{itemize}

% =====================================================================
\section{Floorplan della parte RF}
\label{sec:floorplan}

\subsection{Un vincolo che decide la disposizione}
\begin{attenzione}[ordine dei pin: TX e RX si incrociano se non si gira il Si4464]
I pin dell'HMC349A e del Si4464 sono numerati in senso antiorario (vista dall'alto):
\begin{itemize}
\item IC202, con RFC (pin~3) verso J201: RF2 (pin~8, \textbf{TX}) è in alto a destra, RF1 (pin~5,
      \textbf{RX}) in basso a destra;
\item IC204, con il lato RF verso sinistra: dall'alto RXP (pin~2), RXN (pin~3), poi TX (pin~4).
\end{itemize}
Le due catene formano un anello chiuso dallo switch e dal Si4464. Si dimostra che senza incroci
sullo stesso strato i casi sono due: o la porta RFC dello switch resta \emph{dentro} l'anello (impossibile:
J201 è sul bordo) oppure \textbf{i pin non RF del Si4464 restano dentro l'anello e il suo lato RF guarda
verso l'esterno}. Lo schema del \S24.3 del log (Si4464 con il lato RF verso lo switch, RX sopra e TX
sotto) è indicativo e, così com'è, porta a un incrocio TX/RX vicino allo switch.
\end{attenzione}

\begin{regola}[orientamento di IC202 e IC204]
IC202 con RFC verso J201; \textbf{catena TX sopra} (da RF2) e \textbf{catena RX sotto} (da RF1). IC204
\textbf{ruotato di \ang{180}}, con il lato RF (pin 1--5) verso l'esterno dell'anello: la RX gira intorno
all'angolo inferiore del Si4464 ed entra in RXP/RXN; la TX esce dal pin~4 e gira intorno all'angolo
superiore. Le reti AN643 e AN648 stanno in quei due angoli.
\end{regola}

\begin{regola}[corridoio digitale]
I segnali che devono entrare nell'anello (SPI, nIRQ, SDN di IC204; \net{EN_TX_ANT_SW} per Q701 e l'APC
del PA) lo attraversano \textbf{una sola volta}, a \ang{90}, su B.Cu o In2.Cu, con In1.Cu GND integro
sopra, nel \textbf{tratto TX a basso livello} (tra C701 e C255, sotto l'attenuatore). Mai sotto la catena
RX, la rete AN643, il PA, la rete d'uscita, FL201 o J201.
\end{regola}

\begin{perche}
Il segnale TX tra Si4464 e PA è a \SIrange{0}{20}{dBm} su \SI{50}{\ohm}, con In1 GND nel mezzo:
l'accoppiamento verso una pista su B.Cu è molto basso. Sotto la RX anche un disturbo piccolo arriva al
LNA; sotto il PA ci sono \SI{30}{dBm}. Un'alternativa senza incroci RF è passare sotto il corpo di IC202
(lato GND, pin 6--7 e pad esposto), ma si esce vicino all'antenna: da usare solo se il corridoio TX non è
possibile.
\end{perche}

L'alternativa è scambiare RF1 e RF2 di IC202 e portare CTRL dal segnale invertito di U701; è una modifica
di schema e di logica, fuori dallo scopo di questa guida (Appendice~\ref{app:verifiche}).

\subsection{Disposizione dei blocchi}

\begin{center}
\begin{tikzpicture}[x=1.55mm,y=1.55mm]
% board area
\fill[black!3] (0,0) rectangle (104,64);
\draw[black!40,dash pattern=on 2pt off 1.5pt] (0,0) rectangle (104,64);
\fill[black!70] (0,22) rectangle (0.8,38);
\node[lab,rotate=90] at (-1.8,30) {bordo scheda};
% TX/RX separation fence
\draw[fence] (22,31) -- (74,31);
\node[lab,fill=black!3,inner sep=1pt] at (50,31) {fila di via GND (separazione TX/RX)};
% J201, C206, switch
\node[neu,minimum width=7mm,minimum height=9mm] (j) at (4.6,30) {J201};
\node[neu,minimum width=4mm,minimum height=3mm] (c206) at (10.5,30) {C206};
\node[neu,minimum width=8mm,minimum height=11mm,fill=black!12] (sw) at (17,30) {IC202};
\draw[line width=1.2pt] (j) -- (c206) -- (sw);
% TX row (top, y=48)
\node[txb,minimum height=5mm] (c208) at (24,47) {C208};
\node[txb,minimum width=10mm,minimum height=7mm] (fl) at (31,47) {FL201};
\node[txb] (c230) at (38.5,47) {C230};
\node[pab,minimum width=15mm,minimum height=8mm] (om) at (46.5,47) {C261 C262\\L243 C260};
\node[pab,minimum width=10mm,minimum height=10mm] (pa) at (57,47) {IC2\\PA};
\node[txb,minimum width=11mm] (att) at (68,47) {C255\\R245--248};
\node[txb,minimum width=12mm] (an648) at (82,51) {C701 C704\\L702 C702\\L703};
\node[neu,minimum width=9mm,minimum height=4.5mm,fill=txred!30] (l242) at (47,56.5) {L242};
\node[lab,anchor=south] at (47,59.2) {3V3 PA: pista $\ge$\,1\,mm / pour};
\draw[line width=1.2pt,black!60] (l242.south) -- ($(om.north)+(0.5,0)$);
\draw[line width=1.6pt,black!50] (l242.north) -- (47,59);
% Si4464 rotated, RF side to the right
\node[sib,minimum width=12mm,minimum height=12mm] (si) at (82,31) {IC204\\Si4464\\(ruotato \ang{180})};
\node[lab,anchor=west] at (88.5,34.2) {4 TX};
\node[lab,anchor=west] at (88.5,30.8) {3 RXN};
\node[lab,anchor=west] at (88.5,28.6) {2 RXP};
\node[lab] at (75.5,31) {SPI};
% AN643 corner
\node[rxb,minimum width=12mm] (an643) at (92,17) {C219 L205\\L206 C220};
% TCXO
\node[ckb,minimum width=9mm] (y) at (78,22.5) {Y201\\C204 R201};
% RX row (bottom, y=12)
\node[rxb] (c207) at (24,12) {C207};
\node[rxb,minimum width=12mm] (lpf) at (31.5,12) {C210 L201\\C209};
\node[rxb] (c212) at (39,12) {C212};
\node[rxb,minimum width=10mm,minimum height=8mm] (lna) at (46,12) {Q201\\LNA};
\node[rxb,minimum width=24mm] (f2) at (62,12) {C213 C217 L203\\C216 L204 C218};
% RX arrows
\draw[rxl] (sw.south east) |- (c207.west);
\draw[rxl] (c207) -- (lpf); \draw[rxl] (lpf) -- (c212); \draw[rxl] (c212) -- (lna);
\draw[rxl] (lna) -- (f2);
\draw[rxl] (f2.east) -| (an643.south);
\draw[rxl] (an643.north) |- (88,29.6);
% TX arrows
\draw[txl] (88,34.2) -| (an648.south east);
\draw[txl] (an648.west) -| (att.north);
\draw[txl] (att) -- (pa); \draw[txl] (pa) -- (om); \draw[txl] (om) -- (c230);
\draw[txl] (c230) -- (fl); \draw[txl] (fl) -- (c208);
\draw[txl] (c208.west) -| (sw.north east);
% clock
\draw[draw=olive!80!black,line width=1pt,-{Stealth[length=2mm]}] (y.north) -- (79,25);
% fences around PA and FL201
\draw[fence] (51.2,41) -- (51.2,53); \draw[fence] (62.8,41) -- (62.8,53);
\draw[fence] (31,41.5) -- (31,52.5);
\node[lab,anchor=north] at (31,41.3) {via sotto FL201};
% digital corridor
\draw[dig] (75.6,29) -- (70,29) -- (70,62);
\node[lab,text=violet,anchor=west,align=left] at (70.6,60.5) {corridoio digitale\\su B.Cu/In2, a \ang{90}};
% distances
\draw[dim] (57,41.5) -- (57,13.5);
\node[lab,fill=black!3,inner sep=1pt] at (57,35.5) {$\ge$\,15\,mm};
\draw[dim] (62.5,44) -- (75,24.5);
\node[lab,fill=black!3,inner sep=1pt] at (69.5,37.5) {$\ge$\,15\,mm};
\draw[dim] (6,6) -- (41.5,6);
\node[lab,fill=black!3,inner sep=1pt] at (24,6) {J201 $\to$ Q201: $\le$\,15\,mm di pista};
% legend
\node[lab,anchor=west,text=txred] at (2,61) {\textbf{TX} (fino a \SI{30}{dBm} dopo FL201)};
\node[lab,anchor=west,text=rxblue] at (2,2.5) {\textbf{RX} (piccolo segnale)};
\end{tikzpicture}
\end{center}
\noindent{\footnotesize Floorplan indicativo, non in scala. Le quote marcate sono vincoli; posizione
assoluta, forma della zona RF e lato verso l'MCU li sceglie il progettista. La posizione di L242 sopra la
linea d'uscita deriva dalla EVB ruotata di \ang{180} (\S\ref{sec:pa}).}

\subsection{Ordine di piazzamento}
\begin{enumerate}
\item \textbf{J201 sul bordo} e GND della zona RF (In1 pieno, contorno della zona, fila di via sul bordo).
\item \textbf{IC202 e C206} a ridosso di J201; C207 e C208 sui pin RF1/RF2.
\item \textbf{Rete d'uscita del PA, IC2, L242}: copiata dalla EVB come blocco unico (nota L-01), ruotata
      in modo che l'ingresso guardi verso il Si4464 e l'uscita verso FL201.
\item \textbf{FL201} tra C230 e C208, in linea.
\item \textbf{Q201 con il filtro C210--L201--C209 e C212}, il più vicino possibile a RF1.
\item \textbf{IC204} ruotato, con \textbf{rete AN643} sui pin 2/3 e \textbf{rete AN648} sul pin~4;
      \textbf{Y201} vicino a XIN.
\item Attenuatore e C255 tra AN648 e PA; filtro dopo il LNA tra Q201 e la rete AN643.
\item Disaccoppiamenti (IC204, Y201, Q201, IC202, PA), poi fence, cuciture e corridoio digitale.
\end{enumerate}

\subsection{Distanze critiche}
{\small
\begin{tabularx}{\linewidth}{@{}L{55mm}L{30mm}Y@{}}
\toprule
\textbf{Tra} & \textbf{Valore} & \textbf{Origine}\\
\midrule
J201 -- ingresso Q201 (pista totale) & $\le$\,\SI{15}{mm} & \S\ref{sec:rx}: rumore e captazione; perdita \SI{0.006}{dB/mm}\\
PA (IC2, L242, rete d'uscita) -- Y201 & $\ge$\,\SI{15}{mm} & \S24.3 del log (10--15\,mm); qui si prende il valore alto\\
PA -- catena RX e rete AN643 & $\ge$\,\SI{15}{mm} + fila di via & \S24.3 del log\\
ingresso FL201 -- uscita FL201 (piste) & direzioni opposte, nessun tratto parallelo & \S\ref{sec:fl201}\\
ingresso PA -- uscita PA (piste) & lati opposti del PA, mai affiancate & guadagno \SI{42}{dB}, \S\ref{sec:pa}\\
L242, L243 -- L702, L703, FL201 & $\ge$\,\SI{5}{mm}, assi perpendicolari \dv & nota L-01 regola 8\\
pista digitale su F.Cu -- linea RF & $\ge$\,\SI{2}{mm} con GND cucito in mezzo & \S\ref{sec:generali}\\
C204--R201--XIN (pista totale) & $\le$\,\SI{5}{mm} \dv & \S\ref{sec:tcxo}\\
\bottomrule
\end{tabularx}}

% =====================================================================
\section{Regole per ogni blocco}
\label{sec:blocchi}

% ---------------------------------------------------------------------
\subsection{Connettore MMCX J201 e linea d'antenna}
\label{sec:j201}
J201 è un MMCX da bordo scheda (footprint \path{MMCX_Molex_73415-0961_Horizontal_0.8mm-PCB}):
pad centrale \SI{2.0}{}$\times$\SI{1.1}{mm}, due pad di massa \SI{3.0}{}$\times$\SI{1.4}{mm} ai lati.
\net{RF_ANT} va da J201.1 a C206 (\SI{100}{pF}, DC block dello switch).

\begin{regola}[J201]
\begin{itemize}
\item Pad di massa collegati al GND di F.Cu con almeno 3 via ciascuno, a meno di \SI{0.5}{mm} dal bordo
      del pad; fila di via sul bordo scheda ai due lati del connettore.
\item C206 subito dopo il pad centrale (\S24.5), poi IC202. Pista \net{RF_ANT} dritta, larghezza
      \net{RF_50R}, senza via.
\item Nessuna pista o via di altre reti entro \SI{3}{mm} dal connettore.
\end{itemize}
\end{regola}

\begin{perche}[il pad grande non è un problema a 435 MHz]
Il pad centrale sopra \SI{0.0994}{mm} di prepreg vale circa
$C=\varepsilon_0\cdot4{,}1\cdot(2{,}0\times1{,}1)\,\si{mm^2}/\SI{0.0994}{mm}\approx\SI{0.80}{pF}$, cioè
$|X|=\SI{455}{\ohm}$ a \SI{435}{MHz}: in parallelo a \SI{50}{\ohm} dà $|\Gamma|\approx0{,}05$, return loss
circa \SI{25}{dB}. Non serve togliere il GND di In1 sotto il pad (e non va tolto: è il riferimento).
\end{perche}

\begin{attenzione}[footprint e spessore della scheda]
Il footprint scelto è per un PCB da \SI{0.8}{mm}; il codice JLC C2999837 è un DreamLNK DL-MMCX-KEF-180
(revisione indipendente, P05). Con uno stackup da \SI{1.6}{mm} il connettore edge-mount per \SI{0.8}{mm}
non si monta correttamente: esiste anche \path{MMCX_Molex_73415-0961_Horizontal_1.6mm-PCB}. Brik81
ha chiuso P-05/R-23 senza modifiche; resta da verificare la compatibilità tra codice, footprint e spessore
scelto. \dv
\end{attenzione}

% ---------------------------------------------------------------------
\subsection{Switch HMC349A (IC202)}
\label{sec:sw}
SOP-8 con pad esposto: pin~3 RFC, pin~5 RF1 (RX), pin~8 RF2 (TX), pin 6--7 e pad esposto a GND, pin~1
VDD (5V0, C211), pin~2 CTRL, pin~4 EN a GND. Le porte RF sono in continua verso VDD all'interno, per
questo servono i DC block C206, C207, C208 \fonte{HMC349AMS8G DS rev.~F p.~8}.

\begin{regola}[IC202]
\begin{itemize}
\item \textbf{Pad esposto con matrice di via} (almeno 6, per esempio 2$\times$3 da \SI{0.3}{mm}): la EVB ADI
      mette \enquote{as many plated through vias as possible} attorno alle linee e sotto il pad
      \fonte{DS p.~9}. Il footprint \texttt{SOP65P490X110-9N} oggi \textbf{non ha via} nel pad 9:
      aggiungerle.
\item Pin 6 e 7 collegati al pad esposto e a via proprie; pin~4 (EN) a GND con via sul pad.
\item C206, C207, C208 (0603) a ridosso dei pin RFC, RF1, RF2, in linea con la pista.
\item C211 sul pin~1 con via GND sul pad. CTRL (pin~2) è accanto a RFC: la pista di controllo esce
      perpendicolare, verso l'esterno, senza correre lungo \net{RF_ANT} o \net{RF_SW_COM}.
\end{itemize}
\end{regola}

% ---------------------------------------------------------------------
\subsection{Catena RX: filtro, LNA, filtro post-LNA, rete AN643}
\label{sec:rx}

\subsubsection{Dallo switch al LNA}
RF1 $\rightarrow$ C207 (\SI{100}{pF}) $\rightarrow$ \net{RF_RX_F1} (C210 \SI{6.8}{pF} verso GND)
$\rightarrow$ L201 (\SI{22}{nH}) $\rightarrow$ \net{RF_RX_F2} (C209 \SI{6.8}{pF} verso GND)
$\rightarrow$ C212 (\SI{100}{pF}) $\rightarrow$ Q201 pin~2. È un passa-basso $\pi$ (perdita \SI{0.17}{dB} a
\SI{435}{MHz}, revisione del 04/11, P11; la nota \enquote{passa banda} sullo schema va ignorata, decisione
di Brik81 registrata nel log 043).

\begin{regola}[ingresso LNA]
\begin{itemize}
\item Pista totale da J201 a Q201 pin~2 \textbf{non oltre \SI{15}{mm}}, filtro compreso; componenti in
      linea, senza ritorni.
\item C210 e C209 (0805) con la via di GND sul pad lato massa, non condivisa.
\item L201 (0805) con l'asse lungo la pista; nessun'altra induttanza parallela entro \SI{3}{mm}.
\item Nessuna pista di altre reti tra switch e LNA, su nessuno strato (In1 GND, In2/B.Cu solo GND).
\end{itemize}
\end{regola}

\begin{perche}[perché il LNA va vicino allo switch]
La perdita di \SI{15}{mm} di linea è solo \SI{0.09}{dB} circa (\S\ref{sec:stackup}): conta poco sul rumore.
Il motivo principale è la captazione: davanti al LNA ogni segnale indesiderato raccolto dalla pista viene
amplificato di circa \SI{21}{dB} insieme al segnale utile.
\end{perche}

\subsubsection{LNA TQP3M9036 (Q201)}
Pin~2 RF In, pin~7 RF Out con bias tramite L202 (\SI{120}{nH}, verso 3V3 con C214 \SI{100}{pF} e
C215 \SI{1}{\micro F}), pin~6 Shut Down, pin NC a GND, pad esposto a GND.

\begin{regola}[Q201]
\begin{itemize}
\item \textbf{Pad esposto con via}: il datasheet le richiede per massa RF e calore, foro \SI{0.35}{mm},
      finito \SI{0.25}{mm} \fonte{TQP3M9036 DS p.~11}. Il footprint \texttt{SON50P200X200X95-9N-D} oggi
      \textbf{non ha via}: aggiungerne almeno 3 in colonna (pad \SI{0.85}{}$\times$\SI{1.65}{mm}).
\item Pin NC 1, 3, 4, 5, 8 saldati su pad collegati a GND (\enquote{grounded land pads}, \fonte{DS p.~10}).
\item L202 dal pin~7 con pista corta; C214 sul lato 3V3 di L202 con via sul pad, C215 subito dopo.
      L202 perpendicolare a L203.
\item R202/R203 sul pin~6, lato opposto all'ingresso; la pista di \net{EN_RX_LNA} arriva dall'esterno
      dell'anello senza passare vicino a \net{RF_LNA_IN}.
\item Ingresso (pin~2) e uscita (pin~7) sono su lati opposti: le due piste non devono affiancarsi.
\end{itemize}
\end{regola}

\subsubsection{Filtro dopo il LNA}
C213 (\SI{100}{pF}, serie) $\rightarrow$ \net{RF_RX_F3} (C217 \SI{5.6}{pF}) $\rightarrow$ L203
(\SI{22}{nH}) $\rightarrow$ \net{RF_RX_F4} (C216 \SI{8.2}{pF}) $\rightarrow$ L204 (\SI{22}{nH})
$\rightarrow$ \net{RF_RX_MATCH} (C218 \SI{5.6}{pF}). Passa-basso a 5 elementi.

\begin{regola}[filtro post-LNA]
Componenti in linea, condensatori in parallelo con via propria, L203 e L204 sulla stessa linea ma
\textbf{non affiancate} (distanza tra i corpi $\ge$\,\SI{1}{mm}, assi allineati alla pista). Questo tratto
può essere lungo quanto serve per raggiungere l'angolo inferiore del Si4464.
\end{regola}

\subsubsection{Rete AN643 e coppia RF\_RX\_P / RF\_RX\_N}

\begin{center}
\begin{circuitikz}[scale=0.85,transform shape,font=\sffamily\small]
\draw (0,0) node[left]{\net{RF_RX_MATCH}} to[short,o-*] (1,0);
\draw (1,0) -- (1,1.5) to[C,l=C219 2{,}2p] (4,1.5) to[short,-*] (5,1.5) -- (6.5,1.5) node[right]{RXP (pin 2)};
\draw (1,0) -- (1,-1.5) to[L,l_=L205 56n] (4,-1.5) to[short,-*] (5,-1.5) -- (6.5,-1.5) node[right]{RXN (pin 3)};
\draw (5,1.5) to[L,l=L206 62n] (5,-1.5);
\draw (4,-1.5) to[short,*-] (4,-2.2) to[C,l=C220 4{,}7p] (4,-3.6) node[ground]{};
\draw (1,0) to[C,l_=C218 5{,}6p] (-0.3,-1.6) node[ground]{};
\draw[decorate,decoration={brace,amplitude=4pt,raise=2pt},draw=sbgray] (6.6,1.8) -- (6.6,-1.8)
      node[midway,right=8pt,align=left,font=\sffamily\scriptsize,text=sbgray]{$<$\,3\,mm,\\uguali};
\end{circuitikz}
\end{center}

L'ingresso del LNA del Si4464 a \SI{434}{MHz} vale circa \SI{480}{\ohm} in parallelo a \SI{0.97}{pF}
(AN643 Tab.~1, riportato nel log 007). Non è una linea da adattare in impedenza: conta che le due piste
siano corte, simmetriche e con parassiti uguali (\S24.2).

\begin{regola}[rete AN643]
\begin{itemize}
\item Rete compatta sull'angolo del Si4464 vicino ai pin 2/3; dal nodo \net{RF_RX_MATCH} i due rami
      (C219 verso RXP, L205 verso RXN) partono simmetrici rispetto all'asse tra i due pin.
\item \net{RF_RX_P} e \net{RF_RX_N} più corte di \SI{3}{mm} ciascuna e \textbf{uguali entro
      \SI{0.5}{mm}}; stessa larghezza (\SI{0.20}{mm}), stessa distanza dal GND.
\item L206 (0603) subito tra i pad RXP e RXN, con due tratti di uscita dai pin uguali (i pad sono a
      passo \SI{0.5}{mm}: L206 va di traverso, a meno di \SI{1}{mm} dai pin).
\item C220 (0805) sul lato RXN con la via di GND direttamente sul pad lato massa.
\item Footprint standard (non HandSolder) per C219, C220, L205, L206 (\S24.5). Lasciare \SI{0.5}{mm} libero
      intorno ai quattro componenti per la taratura con VNA prevista da AN643.
\item Nessuna via sulle due reti; nessun'altra pista entro \SI{1}{mm}.
\end{itemize}
\end{regola}

\begin{perche}[i parassiti contano qui]
Il pad di un 0402 vale circa \SI{0.09}{pF}, quello di un 0603 circa \SI{0.2}{pF}
($\varepsilon_0\varepsilon_rA/h$ su \SI{0.0994}{mm}): confrontati con gli \SI{0.97}{pF} dell'ingresso del
Si4464 non sono trascurabili, e se sono diversi tra RXP e RXN sbilanciano la coppia. Una differenza di
lunghezza di \SI{0.5}{mm} invece sposta la fase di appena $360^\circ\cdot0{,}5/388\approx\ang{0.5}$.
\end{perche}

% ---------------------------------------------------------------------
\subsection{Si4464 (IC204)}
\label{sec:si}
QFN-20 \SI{4}{}$\times$\SI{4}{mm}. Pin~2/3 RXP/RXN, pin~4 TX, pin~6 e 8 VDD (3V3, C201 \SI{100}{nF}),
pin~17 XIN, pin~16 XOUT libero, pin~18 GND, pad esposto 21 GND. Il pad esposto fornisce la massa RF e
di tutto il chip: serve una buona saldatura sul piano di massa \fonte{Si4464 DS rev.~1.4 p.~47}.

\begin{regola}[IC204]
\begin{itemize}
\item \textbf{Pad esposto con via} verso In1 (per esempio 3$\times$3 da \SI{0.3}{mm}, passo \SI{0.8}{mm}
      \dv): il footprint \texttt{QFN50P400X400X90-21N-D} oggi \textbf{non ha via}. Pasta a finestre
      (circa 50--70\,\% dell'area \dv) per non far galleggiare il chip.
\item Pin~18 (GND) collegato al pad esposto e a una via propria.
\item C201 a ridosso dei pin~6/8, via GND sul pad. \dv{} Silicon Labs usa nelle sue schede più
      condensatori per VDD (AN629 non consultabile): se il progettista ne aggiunge, il più piccolo va più
      vicino al pin.
\item SPI, nIRQ e SDN escono dal lato opposto al lato RF e raggiungono il corridoio digitale
      (\S\ref{sec:floorplan}); nessuna pista digitale sotto il chip su F.Cu o B.Cu oltre al necessario.
\end{itemize}
\end{regola}

\subsubsection{Rete TX AN648 vicino al pin}
\net{RF_SI_TX}: pin~4, L703 (\SI{220}{nH}, choke da 3V3) e C702 (\SI{9.1}{pF}, serie) $\rightarrow$
\net{RF_TX_M1} $\rightarrow$ L702 (\SI{33}{nH}) $\rightarrow$ \net{RF_TX_M2}: C704 (\SI{15}{pF} verso GND),
C701 (\SI{100}{pF}, DC block verso l'attenuatore). Valori AN648 Tab.~2 per \SI{434}{MHz} (log 001).

\begin{regola}[AN648]
\begin{itemize}
\item L703 e C702 attaccati al pin~4 (pista $\le$\,\SI{1}{mm} \dv); la rete prosegue in linea, lontana
      dalla rete AN643: le due reti partono dal Si4464 in direzioni opposte (TX verso l'angolo superiore,
      RX verso quello inferiore).
\item L703 e L702 con l'asse perpendicolare a L205/L206 (induttanze vicine non allineate: AN629 secondo
      la sintesi di Microwaves\&RF \dv).
\item Lato 3V3 di L703: in netlist non c'è un condensatore dedicato. Portare il 3V3 a L703 con una via vicina
      a C201 o, meglio, prevedere un condensatore (\SI{100}{pF} + \SI{100}{nF}) sul nodo: richiede modifica di
      schema \dv.
\item C704 con via propria sul pad. Se le misure di spurie lo richiedessero, il passa-basso AN648 andrebbe
      tra C704 e C701 (log 001): lasciare, se possibile, circa \SI{4}{mm} di linea dritta in quel punto.
\end{itemize}
\end{regola}

\begin{perche}
Il pin TX lavora in classe E: la tensione di picco sul pin arriva a \SI{6.5}{V} (Si4464 DS tab.~7, citato
nel log 001) e la rete fa parte del funzionamento dello stadio. Il pin~4 è accanto a RXN (pin~3): più
corta è la rete sul pin, meno il TX si accoppia nella rete RX e nel resto della scheda.
\end{perche}

% ---------------------------------------------------------------------
\subsection{TCXO Y201}
\label{sec:tcxo}
ECS-TXO-20CSMV4-320 (\SI{32}{MHz}, \SI{2.0}{}$\times$\SI{1.6}{mm}): pin 1--2 GND, pin~3 uscita, pin~4 VDD
(3V3, C202 \SI{1}{\micro F} e C203 \SI{100}{nF}). Uscita in clipped sine su \SI{10}{k\ohm}\,//\,\SI{10}{pF},
accoppiata in AC: \net{RF_TCXO_OUT} $\rightarrow$ C204 (\SI{10}{nF}) $\rightarrow$ \net{RF_TCXO_AC}
$\rightarrow$ R201 (\SI{33}{\ohm}) $\rightarrow$ \net{RF_XIN} (pin~17, TP701). Il Si4464 chiede
\SIrange{0.6}{1.4}{V} picco-picco in AC su XIN \fonte{Si4464 DS p.~38}.

\begin{regola}[TCXO]
\begin{itemize}
\item Y201 vicino al pin~17, dentro l'anello (lato XIN del Si4464, dopo la rotazione rivolto verso il basso):
      C204 e R201 in linea, pista totale TCXO$\rightarrow$XIN $\le$\,\SI{5}{mm} \dv; netclass \net{RF_CLK}
      (\SI{0.20}{mm}, distanza \SI{0.30}{mm}, \S24.2).
\item C203 e C202 sul pin~4 di Y201, il \SI{100}{nF} più vicino, ciascuno con via GND sul pad.
\item Pin 1--2 a GND con via proprie; GND pieno sotto il TCXO su tutti gli strati.
\item \textbf{Almeno \SI{15}{mm} dal PA, da L242 e dalla pista 3V3 del PA}; fila di via tra il TCXO e la
      catena RX che gli passa sotto.
\item TP701 è un pad da \SI{1}{mm} sulla pista \net{RF_XIN}: va posato direttamente sulla pista, senza
      tratto laterale (vale circa \SI{0.29}{pF}, trascurabile rispetto ai \SI{10}{pF} di carico).
\end{itemize}
\end{regola}

\begin{perche}
Il TCXO non deve ricevere calore (deriva) né RF del PA (spurie e modulazione della frequenza di
riferimento). Le armoniche del clock (13$\times$32 = \SI{416}{MHz}, 14$\times$32 = \SI{448}{MHz}) cadono
vicino alla banda: piste di clock corte e circondate da GND limitano quanto arriva alla catena RX.
\end{perche}

% ---------------------------------------------------------------------
\subsection{Attenuatore R245--R248 e C255}
\label{sec:att}
$\pi$ con R245 \SI{100}{\ohm} (0603) verso GND su \net{RF_ATT_IN}, R246 \SI{68}{\ohm} (0603) in serie,
R247 \SI{100}{\ohm} e R248 \SI{180}{\ohm} (0402) verso GND su \net{RF_ATT_OUT}, poi C255 \SI{56}{pF} verso il
pin~3 del PA. Con carico \SI{50}{\ohm}: $Z_\mathrm{in}=\SI{49.0}{\ohm}$, attenuazione \SI{10.76}{dB}
(calcolo con R247\,//\,R248 = \SI{64.3}{\ohm}). R248 corrisponde a R1 della EVB Qorvo, C255 a C1.

\begin{regola}[attenuatore]
\begin{itemize}
\item Blocco compatto, in linea, ogni resistenza verso GND con via propria sul pad.
\item R248 e C255 copiano l'ingresso della EVB: circa \SI{0.9}{mm} da R248 a C255 e \SI{0.9}{mm} da C255 al
      pin~3 (linee TL2 e TL3 del modello Qorvo, EVB Test Summary p.~12, letti dalla figura \dv).
\item R245 dissipa circa \SI{49}{mW} con il Si4464 a \SI{+20}{dBm} (revisione 04/11): pad 0603 collegato
      a un po' di rame GND.
\item Questo è il punto dove il corridoio digitale può passare sotto (su B.Cu/In2, In1 integro).
\end{itemize}
\end{regola}

% ---------------------------------------------------------------------
\subsection{PA QPA9510 (IC2) e rete d'uscita}
\label{sec:pa}

\subsubsection{Nota L-01 (ripresa dal log, rev.~1.45)}
\begin{tcolorbox}[enhanced,breakable,colback=txred!4,colframe=txred,boxrule=1pt,arc=1.5pt,
  title={\sffamily\bfseries Nota di layout L-01 \textendash{} rete d'uscita del QPA9510},fonttitle=\small,
  coltitle=white,colbacktitle=txred,left=6pt,right=6pt,top=4pt,bottom=4pt]
\small
\begin{enumerate}
\item Copiare la EVB: stesso ordine, stesse distanze, lunghezze e larghezze di pista tra pin RF\_OUT,
      C260, L243 e C261/C262.
\item Se lo stackup è diverso si ricalcola solo la larghezza della linea \SI{50}{\ohm} dopo C230; nel tratto
      pin--C261/C262 restano lunghezze e posizioni della EVB.
\item Tutto sul top, lato PA, nessuna via nel percorso RF, nessun cambio di strato.
\item C260, C261, C262: ognuno con la propria via di massa sul pad o subito accanto (almeno 2 vie per C260
      e C261), non in comune con una pista; piano di massa continuo sotto la rete.
\item Pad esposto del PA: matrice di vie come da datasheet/EVB (massa RF e calore).
\item L242 al nodo \net{RF_PA_OUT} nella stessa posizione della EVB; lato 3V3 disaccoppiato subito con vie sul
      pad; pista larga per circa \SI{1.7}{A} di picco, separata da MCU e radio.
\item Circa \SI{1}{mm} libero intorno a C260 e C262 per sostituirli a mano al banco.
\item L243 e L242 lontane dalle induttanze del match TX del Si4464 e da FL201.
\end{enumerate}
\end{tcolorbox}

\subsubsection{Corrispondenza con la EVB Qorvo 410--470 MHz}
{\footnotesize
\begin{tabularx}{\linewidth}{@{}lllY@{}}
\toprule
\textbf{SYMPHONY} & \textbf{EVB} & \textbf{Valore} & \textbf{Posizione / funzione}\\
\midrule
C255 & C1 & 56p & DC block ingresso, a circa \SI{0.9}{mm} dal pin~3\\
R248 & R1 & 180 & resistenza verso GND sull'ingresso (prima di C1)\\
L240 / R240 & L1 / L7 & 12n / 0\,\si{\ohm} & VCC1, match interstadio; L7 è DNP sulla EVB, qui R240 montata\\
C244 / C240 / C241 & C19 / C2 / C3 & 330p / 10n / 1n & disaccoppiamento VCC1\\
L241 / R241 & L6 / L2 & 6,8n / 0\,\si{\ohm} & VCC2 (pin 5, 6, 7), match interstadio\\
C248 / C245 / C246 / C247 & C20 / C21 / C5 / C6 & 330p / 10\textmu{} (3,3\textmu) / 10n / 1n & disaccoppiamento VCC2\\
C251 & C15 & 100p & VCC pin~14\\
C252 & C16 & 10n & APC (pin 15/16); C17 della EVB non ha corrispondente\\
L242 & L3 & 1\si{\micro H} & choke dello stadio finale (NRS2012, 2020)\\
C256 / C257 / C258 & C18 / C13 / C14 & 10\textmu{} (3,3\textmu) / 1n / 47p & disaccoppiamento lato 3V3 di L242\\
C260 & C10 & 27p (0603) & primo condensatore in parallelo al nodo \net{RF_PA_OUT}\\
L243 & L4 & 2,7n (0402) & serie\\
C261 / C262 & C11 / C25 & 22p / 2p (0402) & in parallelo al nodo \net{RF_PA_MATCH} (C262 = taratura)\\
C230 & C12 & 56p & DC block verso la linea \SI{50}{\ohm}\\
\bottomrule
\end{tabularx}}
\noindent{\footnotesize Fonti: EVB Test Summary p.~2--3 (schema, foto, BOM), revisione 04/11 \S2.2. C8 (2F0), C9,
C22--C24 della EVB sono DNP.}

\subsubsection{Geometria della EVB e cosa cambia con lo stackup}
La pagina~12 della EVB Test Summary mostra il modello usato da Qorvo per l'impedenza di carico. Lette
dalla figura (risoluzione bassa \dv), le linee d'uscita sono:

\begin{center}
\begin{circuitikz}[scale=0.74,transform shape,font=\sffamily\small]
\draw (0,0) node[left,align=right]{pin RF\_OUT\\(9--12)} to[TL,l=1{,}34] (2.6,0) to[short,-*] (2.6,0);
\draw (2.6,0) to[C,l_=C260 27p] (2.6,-1.8) node[ground]{};
\draw (2.6,0) -- (2.6,0.9) to[L,l=L242 1µ] (2.6,2.6) node[above]{3V3};
\draw (2.6,0) to[TL,l=1{,}54] (5.2,0) to[L,l=L243 2{,}7n] (7.4,0) to[TL,l=0{,}88] (9.6,0) to[short,-*] (9.6,0);
\draw (9.6,0) to[C,l_=C261 22p] (9.6,-1.8) node[ground]{};
\draw (9.6,0) -- (10.8,0) to[short,-*] (10.8,0) to[C,l=C262 2p] (10.8,-1.8) node[ground]{};
\draw (10.8,0) to[TL,l=3{,}82] (13.4,0) to[C,l=C230 56p] (15.4,0) to[TL,l=5{,}97] (17.8,0) node[right]{FL201};
\end{circuitikz}
\end{center}
\noindent{\footnotesize Lunghezze in mm tra i componenti, larghezza del modello \enquote{W = 19,8}
(probabilmente mil, circa \SI{0.50}{mm}). Il punto d'innesto di L242 (L3) va rilevato dalla foto/Gerber EVB.}

\begin{perche}[perché le piste fanno parte del match]
Al pin l'obiettivo Qorvo è $4{,}47+j0{,}65\,\si{\ohm}$ a \SI{430}{MHz} \fonte{EVB Test Summary p.~12}. Una
linea con $Z_0$ e $\varepsilon_\mathrm{eff}$ ha induttanza $L'=Z_0\sqrt{\varepsilon_\mathrm{eff}}/c$: per una
linea da \SI{50}{\ohm} ($\varepsilon_\mathrm{eff}=3{,}16$) $L'=\SI{0.30}{nH/mm}$, cioè circa
\SI{0.8}{\ohm} di reattanza per millimetro a \SI{435}{MHz}, una frazione grande dei \SI{4.5}{\ohm} del
nodo. Calcoli fatti per questa guida (ABCD, linee senza perdite, $\varepsilon_\mathrm{eff}=2{,}8$):
\begin{itemize}
\item componenti ideali, senza linee: $Z=1{,}78-j4{,}60\,\si{\ohm}$ a \SI{435}{MHz} (stesso valore della
      revisione del 04/11);
\item con le lunghezze lette dalla figura: $Z=2{,}16-j2{,}68\,\si{\ohm}$ (\SI{435}{MHz}); con ESL tipiche
      dei condensatori (\SIrange{0.4}{0.6}{nH} \dv) $Z\approx2{,}0-j2{,}2\,\si{\ohm}$.
\end{itemize}
Le linee spostano l'impedenza verso l'obiettivo ma, con i dati disponibili, non bastano a raggiungerlo:
il modello Qorvo contiene anche pad, gradini e modelli dei componenti che dalla figura non si leggono.
Conclusione: copiare la EVB il più fedelmente possibile e tarare al banco (M-06), come dice L-01.
\end{perche}

\begin{attenzione}[L-01 punto 1 e 2 con uno stackup diverso dalla EVB]
Le piste della EVB (circa \SI{0.5}{mm}, compatibili con \SI{50}{\ohm} su \SI{0.25}{mm} di RO4350B, lo stesso
substrato dichiarato per la EVB della TQP3M9036 \fonte{TQP3M9036 DS p.~10}; substrato della EVB QPA9510
\dv) sul JLC04161H-3313 valgono circa \SI{26}{\ohm}: a parità di lunghezza hanno circa metà
dell'induttanza ($L'\approx\SI{0.15}{nH/mm}$ contro \SI{0.28}{nH/mm}). Copiare la larghezza fisica
(L-01 punto~1) non copia quindi il comportamento elettrico. Questa guida segue L-01 alla lettera
(stesse posizioni, lunghezze e larghezze, \SI{0.5}{mm}), perché la pista deve comunque portare la corrente
(vedi sotto), e conta sulla taratura di C260 e C262. Se dopo la taratura la potenza resta bassa, la
prima cosa da rivedere è questo tratto. Da decidere con Brik81 (Appendice~\ref{app:verifiche}).
\end{attenzione}

\subsubsection{Regole per il PA}
\begin{regola}[rete d'uscita, integra L-01]
\begin{itemize}
\item \textbf{Blocco EVB ruotato di \ang{180}}: ingresso (pin~3) verso l'attenuatore, uscita (pin 9--12)
      verso FL201. La rotazione conserva tutte le distanze; il lato APC (pin 13--16) guarda dentro l'anello,
      il lato VCC2 (pin 5--8) e L242 verso l'esterno.
\item Pin 9--12 uniti da un unico pad di rame sul lato del package (come sulla EVB), da cui parte la pista
      d'uscita.
\item Pista pin$\rightarrow$C260$\rightarrow$L243$\rightarrow$C261/C262 larga come sulla EVB (circa
      \SI{0.5}{mm} \dv), lunghezze della figura sopra; dopo C262 la linea passa a \net{RF_50R}
      (\SI{0.20}{mm}) con una rastremazione di circa \SI{0.5}{mm}.
\item Masse dirette (\enquote{solid}, senza razze termiche) per C260, C261, C262, C230: almeno 2 via per
      C260 e C261, 1--2 per C262, sul pad o a meno di \SI{0.3}{mm}.
\item \SI{1}{mm} libero intorno a C260 e C262; niente serigrafia sopra i pad; riferimenti leggibili.
\item Pin~8 (2F0) e pin~13 (NC): pad saldati ma \textbf{senza pista}; il pin~7 (NC) va unito a VCC2 come
      sulla EVB (decisione P19).
\end{itemize}
\end{regola}

\begin{regola}[masse del PA]
\begin{itemize}
\item \textbf{Pin~2 (GND1)}: collegamento corto, via dedicata subito accanto
      \fonte{QPA9510 DS p.~7}.
\item \textbf{Pin~4 (GND2)}: il datasheet consiglia una pista di circa \textbf{40 mil (\SI{1}{mm})} verso il
      piano di massa, per ridurre il rumore in uscita (riduce di poco il guadagno), e via dedicate per la
      stabilità \fonte{QPA9510 DS p.~7}. Copiare la forma della EVB.
\item Disaccoppiamenti: il condensatore più piccolo (\SI{47}{}/\SI{100}{}/\SI{330}{pF}) attaccato al pin o
      al nodo, poi \SI{1}{nF}, poi \SI{10}{nF}, poi \SI{10}{\micro F}; ogni condensatore con la sua via
      (\S24.4). L240 e L241 nella posizione della EVB (fanno parte del match interstadio,
      \fonte{QPA9510 DS p.~7}).
\end{itemize}
\end{regola}

\begin{regola}[pad termico e calore]
\begin{itemize}
\item Matrice 3$\times$3 nel pad esposto (già nel footprint: foro \SI{0.25}{mm}, pad \SI{0.45}{mm}, passo
      \SI{0.5}{mm}); Qorvo indica foro finito \SIrange{0.203}{0.330}{mm} su griglia \SIrange{0.5}{1.2}{mm} e
      rame da 1 oz su top e bottom \fonte{QPA9510 DS p.~8}.
\item Aggiungere \textbf{almeno 12 via} nel GND di F.Cu intorno al package (entro \SI{3}{mm}) e un'area di
      rame su B.Cu di almeno \SI{10}{}$\times$\SI{10}{mm} \dv, accoppiabile termicamente alla struttura
      del CubeSat.
\item Via del pad riempite e chiuse se JLC lo offre sul 4 strati (sulla pagina capacità è predefinito
      dal 6 strati in su \dv); altrimenti chiuse dal lato B.Cu con solder mask e pasta a finestre sul pad.
\end{itemize}
\end{regola}

\begin{perche}[il conto termico]
Il PA dissipa circa \SIrange{2.1}{2.3}{W} alla massima potenza (\S24.4 e log 001). Una via con foro
\SI{0.25}{mm} e rame del foro di \SI{20}{\micro m} (valore assunto \dv) attraverso \SI{1.6}{mm} ha
$R_\theta=\ell/(k\,\pi d\,t)=\num{1.6e-3}/(385\cdot\pi\cdot\num{0.25e-3}\cdot\num{20e-6})\approx\SI{265}{K/W}$;
9 via in parallelo $\approx\SI{29}{K/W}$, cioè circa \SI{65}{K} con \SI{2.2}{W} se tutto il calore passasse
solo da lì. Per questo servono le via aggiuntive, il rame di B.Cu e il collegamento alla struttura (nel
vuoto non c'è convezione). Con le trasmissioni lunghe verificare la temperatura (S-14, F-03).
\end{perche}

\subsubsection{Alimentazione del PA e L242}
\begin{regola}[3V3 del PA]
\begin{itemize}
\item Pista \textbf{dedicata} dal buck al PA, separata da quelle di MCU, Si4464 e TCXO (log 001).
\item Ramo 3V3$\rightarrow$L242: \textbf{almeno \SI{0.62}{mm} su strato esterno}, consigliato \SI{1}{mm} o
      un pour (\S24.4). Su strato interno JLC (\SI{0.5}{oz}) servirebbero circa \SI{3.7}{mm}: tenere il
      ramo su F.Cu o usare un pour largo.
\item Almeno 4 via da \SI{0.4}{mm} per ogni cambio di strato del ramo di potenza (\S24.4).
\item C258 (\SI{47}{pF}) attaccato al lato 3V3 di L242, poi C257 (\SI{1}{nF}) e C256 (\SI{10}{\micro F}),
      ognuno con via sul pad.
\item L242 lontana da L243, da L702/L703 e da FL201 (L-01 punto 8), corpo dell'induttore non sopra linee RF.
\end{itemize}
\end{regola}

\begin{perche}[i numeri della corrente]
IPC-2221: $I=k\,\Delta T^{0{,}44}A^{0{,}725}$ ($A$ in mil$^2$, $k=0{,}048$ esterno, $0{,}024$ interno). Con
$I=\SI{1.7}{A}$ e $\Delta T=\SI{10}{K}$: esterno $A=\SI{33.9}{mil^2}$, cioè \SI{0.62}{mm} su
\SI{35}{\micro m}; interno $A=\SI{88}{mil^2}$, cioè \SI{3.7}{mm} su \SI{15}{\micro m} (il \S24.4 indica
\SI{1.6}{mm}, valore che vale per 1 oz interno). Caduta su \SI{30}{mm} di pista da \SI{1}{mm}:
\SI{14.7}{m\ohm}, \SI{25}{mV} a \SI{1.7}{A}. Nel tratto pin--C260 passano anche circa \SI{0.67}{A} RF efficaci
($\sqrt{\SI{2}{W}/\SI{4.5}{\ohm}}$): è un altro motivo per non stringere quella pista sotto la larghezza
della EVB.
\end{perche}

\begin{regola}[ingresso e isolamento del PA]
Ingresso e uscita del PA su lati opposti, piste mai affiancate, fila di via GND che separa il tratto
d'ingresso (attenuatore, C255) dalla rete d'uscita e da FL201.
\end{regola}

\begin{perche}
Il PA ha circa \SI{42}{dB} di guadagno a \SI{430}{MHz} \fonte{EVB Test Summary p.~6}: se l'uscita torna
all'ingresso con meno di \SI{42}{dB} di isolamento più margine, il PA oscilla. L'accoppiamento tra piste
vicine su FR4 sottile è dell'ordine di $1/(1+(d/h)^2)$: a $d=\SI{2}{mm}$ e $h=\SI{0.1}{mm}$ circa
\SI{-52}{dB}, appena sufficiente; per questo, oltre alla distanza, serve la fila di via.
\end{perche}

% ---------------------------------------------------------------------
\subsection{Filtro FL201 (LFCG-490+)}
\label{sec:fl201}
Mini-Circuits LFCG-490+, \SI{2.0}{}$\times$\SI{1.25}{mm}: pad~8 ingresso (\net{RF_LPF_IN}, da C230), pad~4
uscita (\net{RF_LPF_OUT}, verso C208), pad 1, 2, 3, 5, 6, 7 GND \fonte{LFCG-490+ DS p.~3}. Perdita circa
\SI{1.15}{dB} a \SI{435}{MHz}; \SI{3.5}{W} a \SI{25}{\celsius} con derating fino a \SI{0.9}{W} a
\SI{125}{\celsius} \fonte{DS p.~1--2}.

\begin{center}
\begin{tikzpicture}[x=6mm,y=6mm]
\fill[black!6] (-6,-2.6) rectangle (6,2.6);
\fill[orange!70!black] (-6,-0.17) rectangle (-1.7,0.17);
\fill[orange!70!black] (1.7,-0.17) rectangle (6,0.17);
\node[font=\sffamily\tiny,above] at (-4.2,0.2) {da C230};
\node[font=\sffamily\tiny,above] at (4.2,0.2) {verso C208};
\draw[black!70,line width=0.5pt,fill=black!20] (-1.0,-0.62) rectangle (1.0,0.62);
\node[font=\sffamily\tiny] at (0,0) {FL201};
\foreach \x in {-0.65,0,0.65}{
  \fill[orange!70!black] (\x-0.18,0.4) rectangle (\x+0.18,1.4);
  \fill[orange!70!black] (\x-0.18,-1.4) rectangle (\x+0.18,-0.4);
  \fill[white] (\x,1.75) circle (0.15); \draw (\x,1.75) circle (0.15);
  \fill[white] (\x,-1.75) circle (0.15); \draw (\x,-1.75) circle (0.15);
}
\fill[orange!70!black] (-1.7,-0.18) rectangle (-0.7,0.18);
\fill[orange!70!black] (0.7,-0.18) rectangle (1.7,0.18);
\foreach \y in {-2.3,2.3} \foreach \x in {-5,-4,-3,-2,2,3,4,5} {\draw[fill=white] (\x,\y) circle (0.12);}
\foreach \y in {-1.0,1.0} \foreach \x in {-5,-4,-3,-2.2,2.2,3,4,5} {\draw[fill=white] (\x,\y) circle (0.12);}
\node[font=\sffamily\tiny,align=center] at (0,-3.2) {via su ogni pad GND, fence ai lati delle linee, nessun tratto parallelo tra ingresso e uscita};
\end{tikzpicture}
\end{center}

\begin{regola}[FL201]
\begin{itemize}
\item In linea tra C230 e C208: ingresso e uscita \textbf{in direzioni opposte}, piste dritte per almeno
      \SI{2}{mm} da ogni lato del filtro; mai un percorso a U che riavvicini l'uscita all'ingresso.
\item Una via per ciascuno dei 6 pad GND, il più vicino possibile al pad (in totale 6 via), più la fence
      delle linee; le via sui due lati del filtro formano una barriera trasversale tra il lato PA e il lato
      switch.
\item Footprint conforme al layout consigliato Mini-Circuits PL-429 (\fonte{DS p.~3}): il footprint
      \texttt{IC:LFCG490} (pad \SI{0.35}{}$\times$\SI{1.0}{mm}) va confrontato con PL-429 \dv.
\item Lontano da L242 e L243 (L-01 punto 8). Il filtro scalda: GND attorno con buone via verso In1.
\end{itemize}
\end{regola}

\begin{perche}
FL201 attenua circa \SI{60}{dB} a 800--960\,MHz \fonte{revisione 04/11, \S2.2}: perché questa attenuazione
arrivi all'antenna, l'accoppiamento diretto tra le piste d'ingresso e d'uscita deve essere molto più
basso di \SI{-60}{dB}. Con il filtro lungo appena \SI{2}{mm} questo si ottiene solo con piste che si
allontanano in direzioni opposte e con le masse sui pad.
\end{perche}

\begin{attenzione}[armoniche dello switch]
FL201 è a monte dello switch: le armoniche prodotte da IC202 a circa \SI{30}{dBm} non vengono filtrate
(revisione 04/11 P10, log M-05). Il layout non lo risolve; tenere solo il tratto C208--RF2 corto.
\end{attenzione}

% =====================================================================
\section{Regole generali RF}
\label{sec:generali}

\subsection{Piano di massa e ritorni di corrente}
\begin{regola}[GND]
\begin{itemize}
\item In1.Cu GND pieno sotto tutta la zona RF (e, se possibile, sotto tutta la scheda); nessuna pista su In1.
\item Sotto la zona RF anche In2.Cu e B.Cu solo GND, cuciti con In1, tranne il corridoio digitale.
\item GND di F.Cu riempito intorno a tutte le linee e i componenti RF, con la distanza della CPWG
      (\SI{0.20}{mm}); nessuna isola di rame scollegata o \enquote{appesa} con una sola via.
\item Ogni condensatore in parallelo ha la sua via: il ritorno non deve passare per la pista o per un pad
      di un altro componente.
\item Le correnti forti del PA (ritorno della pista 3V3 e del buck) non devono attraversare la zona RX:
      la pista 3V3 del PA arriva dall'esterno dell'anello, lato L242.
\end{itemize}
\end{regola}

\subsection{Via di cucitura e via fence}
Il passo delle via segue la regola pratica $p\le\lambda/20$ alla frequenza più alta da contenere (regola
d'uso comune, non da una application note specifica). Con $\lambda_g=\lambda_0/\sqrt{\varepsilon_\mathrm{eff}}$
($\varepsilon_\mathrm{eff}=3{,}16$) e, più conservativo, $\lambda_d=\lambda_0/\sqrt{\varepsilon_r}$ ($\varepsilon_r=4{,}1$):

\begin{center}
\begin{minipage}{0.47\linewidth}
\small
\begin{tabular}{@{}rrrr@{}}
\toprule
$n$ & $f$ [MHz] & $\lambda_g/20$ [mm] & $\lambda_d/20$ [mm]\\
\midrule
1 & 435 & 19,4 & 17,0\\
2 & 870 & 9,7 & 8,5\\
3 & 1305 & 6,5 & 5,7\\
5 & 2175 & 3,9 & 3,4\\
10 & 4350 & 1,9 & 1,7\\
\bottomrule
\end{tabular}
\end{minipage}\hfill
\begin{minipage}{0.5\linewidth}
\begin{tikzpicture}
\begin{axis}[width=\linewidth,height=44mm,xlabel={armonica $n$},ylabel={passo max [mm]},
  ymode=log,xmin=1,xmax=10,ymin=1,ymax=25,grid=major,
  tick label style={font=\scriptsize},label style={font=\scriptsize},
  legend style={font=\tiny,at={(0.98,0.98)},anchor=north east},
  ytick={1,2,5,10,20},yticklabels={1,2,5,10,20}]
\addplot[domain=1:10,samples=10,thick,sbblue] {388/(20*x)};
\addplot[domain=1:10,samples=10,thick,dashed,txred] {340.6/(20*x)};
\addplot[black!60,thin] coordinates {(1,2) (10,2)};
\legend{$\lambda_g/20$,$\lambda_d/20$,2\,mm}
\end{axis}
\end{tikzpicture}
\end{minipage}
\end{center}

\begin{regola}[passo delle via]
\begin{itemize}
\item \textbf{Fence lungo le linee RF}: via da \SI{0.6}{}/\SI{0.3}{mm} su entrambi i lati, centro a
      \SIrange{0.5}{0.7}{mm} dal bordo della pista (cioè a \SIrange{0.3}{0.5}{mm} dal bordo del GND
      coplanare), \textbf{passo $\le$\,\SI{2}{mm}} (\S24.1): copre fino alla 10$^a$ armonica
      (\SI{4.35}{GHz}).
\item \textbf{Intorno al PA, alla rete d'uscita e a FL201: passo $\le$\,\SI{1.5}{mm}} (sotto $\lambda_d/20$ a
      \SI{4.35}{GHz}).
\item Bordo della zona RF e bordo scheda: fila di via a passo $\le$\,\SI{2}{mm}.
\item Cucitura delle aree GND nella zona RF: griglia $\le$\,\SI{3}{mm} (sotto $\lambda/20$ alla 5$^a$
      armonica).
\item Fila di via tra zona TX e zona RX (\S24.3), passo $\le$\,\SI{2}{mm}.
\end{itemize}
\end{regola}

\begin{center}
\begin{tikzpicture}[x=5mm,y=5mm]
\fill[black!8] (0,0) rectangle (16,5);
\fill[white] (0,2.0) rectangle (16,3.0);
\fill[orange!70!black] (0,2.2) rectangle (16,2.8);
\foreach \x in {1,3,...,15} {\draw[fill=white] (\x,3.65) circle (0.22); \draw[fill=white] (\x,1.35) circle (0.22);}
\draw[dim] (1,4.5) -- (3,4.5); \node[lab,above] at (2,4.5) {$\le$\,2\,mm};
\draw[dim] (16.6,2.2) -- (16.6,2.8); \node[lab,right] at (16.7,2.5) {$W=0{,}20$};
\draw[dim] (17.9,2.8) -- (17.9,3.0); \node[lab,right] at (18.0,3.3) {$s=0{,}20$};
\node[lab,below] at (8,0) {CPWG su F.Cu, In1.Cu GND sotto; via di fence sui due bordi del GND coplanare};
\end{tikzpicture}
\end{center}

\subsection{Clearance, angoli, pad e transizioni}
\begin{regola}[geometria delle piste RF]
\begin{itemize}
\item Clearance delle reti RF verso il GND = gap della CPWG (\SI{0.20}{mm}); verso altre reti
      $\ge$\,\SI{0.5}{mm}, verso piste digitali su F.Cu $\ge$\,\SI{2}{mm} con GND cucito in mezzo.
\item Curve a \ang{45} o ad arco (\S24.1). A \SI{435}{MHz} un angolo a \ang{90} su \SI{0.2}{mm} vale circa
      $W\cdot C'\approx\SI{0.02}{pF}$ in più ($C'=\SI{0.12}{pF/mm}$): trascurabile sulle linee \SI{50}{\ohm};
      la regola serve per pulizia e armoniche. Nella rete del PA vale la geometria EVB.
\item Pista che entra in un pad: al centro del lato corto, in linea; nessun pad attraversato di lato.
      La pista non si allarga prima del pad (niente teardrop sulle reti RF).
\item Nessuno stub: la pista va da pad a pad; un condensatore in parallelo si posa \emph{sulla} pista, non
      a lato con un tratto di collegamento.
\item Componenti RF 0402/0603/0805 con footprint standard (non HandSolder) nel percorso RF (\S24.5,
      rev.~1.13). Restano HandSolder solo disaccoppiamenti del PA, L240, L241, L202, R240--R243: accettabile
      perché non sono nel percorso a \SI{50}{\ohm} \dv{} (L240 e L241 fanno parte del match interstadio).
\item Nessuna via nel percorso RF; se inevitabile (non sulle reti di L-01), 2--4 via GND accanto (\S24.5).
\end{itemize}
\end{regola}

\subsection{Separazione dal digitale}
\begin{regola}[digitale]
\begin{itemize}
\item Nessuna pista digitale o di alimentazione switching sotto le linee RF, su nessuno strato, tranne
      l'attraversamento unico del corridoio (\S\ref{sec:floorplan}).
\item SPI, nIRQ, SDN, \net{EN_TX_ANT_SW}, \net{EN_RX_LNA}: dal lato MCU, con il percorso più corto, mai in
      parallelo a una linea RF per più di \SI{5}{mm}.
\item Il buck TPSM863257 e le sue piste di commutazione lontani dalla zona RF (almeno \SI{20}{mm} \dv) e
      mai sotto.
\end{itemize}
\end{regola}

\begin{perche}
Per due piste sopra un piano a distanza $h$ l'accoppiamento scende circa come $1/(1+(d/h)^2)$: con
$h=\SI{0.1}{mm}$, a $d=\SI{1}{mm}$ vale circa \SI{-40}{dB}, a \SI{2}{mm} circa \SI{-52}{dB}. Con un GND
cucito in mezzo il valore reale è ancora più basso.
\end{perche}

\subsection{Schermatura (opzione)}
Lo schermo non è in schema. Se si vuole tenere l'opzione:
\begin{itemize}
\item solo telaio o clip SMD, niente componenti a foro passante (\S24.3);
\item anello di GND su F.Cu senza solder mask, largo circa \SI{1}{mm} \dv, con via a passo
      $\le$\,\SI{2}{mm}; due scomparti: (1) PA, rete d'uscita, FL201; (2) LNA, filtri, Si4464, TCXO;
      lo switch e J201 fuori o in uno scomparto a parte;
\item dove una linea RF attraversa una parete, la parete ha una tacca larga almeno 3 volte $(W+2s)$ \dv;
\item altezza interna maggiore del componente più alto più \SI{0.5}{mm} (L242 \SI{1.2}{mm});
\item richiede l'aggiunta del componente in schema e in BOM (modifica da approvare).
\end{itemize}

\subsection{Punti di test e taratura}
\begin{itemize}
\item Esistenti: TP701 su \net{RF_XIN} (ampiezza TCXO, \SIrange{0.6}{1.4}{Vpp}), TP702 su
      \net{EN_TX_ANT_SW}, TP703 su \net{EN_RX_LNA}. Pad sulla pista, senza tratti laterali.
\item Misure al collaudo dal connettore J201 con attenuatore di potenza (M-01, M-05, M-06 del log).
\item \textbf{Senza modifiche di schema}: lasciare piazzole GND senza solder mask (circa
      \SI{1.5}{}$\times$\SI{1.5}{mm}, su F.Mask) accanto a C230, C255, C212 e C213, per saldare la calza di un
      cavo semirigido al posto del condensatore serie e misurare il PA o il LNA separatamente.
\item \textbf{Opzione con modifica di schema}: un connettore U.FL DNP con 0402 da \SI{0}{\ohm} a T
      sull'uscita di C230 \dv. Non consigliato nel tratto L-01 (aggiunge parassiti al nodo).
\item Spazio di taratura: \SI{1}{mm} intorno a C260 e C262; \SI{0.5}{mm} intorno a C219, C220, L205, L206,
      C702, C704, L702.
\end{itemize}

% =====================================================================
\section{Netclass e regole DRC in KiCad}
\label{sec:netclass}

\subsection{Cosa c'è oggi in SYMPHONY.kicad\_pro}
\begin{attenzione}[netclass mancanti]
Il file contiene una sola classe, \texttt{Default} (pista \SI{0.2}{mm}, clearance \SI{0.2}{mm}, via
\SI{0.6}{}/\SI{0.3}{mm}), con \texttt{netclass\_assignments} vuoto e \texttt{netclass\_patterns} vuoto; la
netlist assegna \texttt{Default} a tutte le reti. Il log dice che in rev.~1.12 sono state aggiunte 6 netclass
con 31 assegnazioni e che in rev.~1.13 le assegnazioni sono state riportate sui nuovi nomi; nessuno dei
backup in \texttt{report\_src/bak\_*} contiene \texttt{RF\_50R}. \textbf{Tutte le netclass e le assegnazioni
vanno ricreate.}
\end{attenzione}

Regole di scheda presenti: clearance minima \SI{0.15}{mm}, pista minima \SI{0.1}{mm}, via minima
\SI{0.3}{mm}, foro passante minimo \SI{0.3}{mm}, anello minimo \SI{0.1}{mm}, foro--foro \SI{0.2}{mm},
rame--bordo \SI{0.2}{mm}.

\begin{attenzione}[via del pad di IC2]
Il footprint di IC2 ha via con foro \SI{0.25}{mm}, sotto il minimo di scheda (\texttt{min\_through\_hole\_diameter}
= \SI{0.3}{mm}): il DRC le segnalerà. JLCPCB accetta fori da \SI{0.2}{mm} (preferito) e \SI{0.15}{mm}
(minimo) sul 4 strati; con foro \SI{0.25}{mm} e via $\ge$\,\SI{0.45}{mm} non c'è sovrapprezzo
\fonte{jlcpcb.com/capabilities}. Abbassare il minimo a \SI{0.25}{mm} o portare le via a \SI{0.3}{mm}
(Qorvo ammette fino a \SI{0.330}{mm}).
\end{attenzione}

\subsection{Netclass da creare}
{\small
\renewcommand{\arraystretch}{1.15}
\begin{xltabular}{\linewidth}{@{}L{22mm}L{23mm}L{17mm}Y@{}}
\toprule
\textbf{Netclass} & \textbf{Pista / clearance} & \textbf{Via} & \textbf{Reti (nomi netlist, prefisso \net{/UHF/})}\\
\midrule\endhead
\midrule\multicolumn{4}{r}{\footnotesize\itshape continua\ldots}\\\endfoot
\bottomrule\endlastfoot
\textbf{RF\_50R} & \SI{0.20}{} / \SI{0.20}{mm} & \SI{0.6}{} / \SI{0.3}{} & \net{RF_ANT} \net{RF_SW_COM} \net{RF_SW_RX} \net{RF_SW_TX} \net{RF_RX_F1} \net{RF_RX_F2} \net{RF_LNA_IN} \net{RF_LNA_OUT} \net{RF_RX_F3} \net{RF_RX_F4} \net{RF_RX_MATCH} \net{RF_SI_TX} \net{RF_TX_M1} \net{RF_TX_M2} \net{RF_ATT_IN} \net{RF_ATT_OUT} \net{RF_PA_IN} \net{RF_PA_MATCH} \net{RF_LPF_IN} \net{RF_LPF_OUT} (20 reti, come \S24.2)\\
\midrule
\textbf{RF\_DIFF} & \SI{0.20}{mm}, gap \SI{0.20}{mm} & -- & \net{RF_RX_P}, \net{RF_RX_N} (coppia: KiCad riconosce i suffissi \_P/\_N)\\
\midrule
\textbf{RF\_PA\_OUT} & \SI{1.0}{} / \SI{0.25}{mm} & \SI{0.8}{} / \SI{0.4}{} & \net{RF_PA_OUT}. Il valore \SI{1.0}{mm} (\S24.2) vale per il ramo L242; il tratto pin--C260--L243 si traccia a mano con la larghezza EVB (\S\ref{sec:pa})\\
\midrule
\textbf{RF\_CLK} & \SI{0.20}{} / \SI{0.30}{mm} & -- & \net{RF_TCXO_OUT}, \net{RF_TCXO_AC}, \net{RF_XIN}\\
\midrule
\textbf{PA\_BIAS} & \SI{0.50}{} / \SI{0.20}{mm} & \SI{0.8}{} / \SI{0.4}{} & \net{PA_VCC1}, \net{PA_VCC1_F}, \net{PA_VCC2}, \net{PA_VCC2_F}\\
\midrule
\textbf{PWR} & \SI{0.50}{} / \SI{0.20}{mm} & \SI{0.8}{} / \SI{0.4}{} & \net{3V3}, \net{5V0} (reti globali, senza prefisso). Il ramo 3V3 del PA si traccia a mano $\ge$\,\SI{1}{mm} o a pour\\
\end{xltabular}}

Pattern consigliati (Impostazioni scheda $\rightarrow$ Netclass $\rightarrow$ assegnazioni), in quest'ordine:
\begin{center}
\small
\begin{tabular}{@{}ll@{}}
\toprule
\textbf{Pattern} & \textbf{Netclass}\\
\midrule
\texttt{/UHF/RF\_RX\_P}, \texttt{/UHF/RF\_RX\_N} & RF\_DIFF\\
\texttt{/UHF/RF\_TCXO\_*}, \texttt{/UHF/RF\_XIN} & RF\_CLK\\
\texttt{/UHF/RF\_PA\_OUT} & RF\_PA\_OUT\\
\texttt{/UHF/PA\_VCC*} & PA\_BIAS\\
le 20 reti di RF\_50R, una per riga (non usare \texttt{/UHF/RF\_*}: prenderebbe anche \net{RF_SDN}) & RF\_50R\\
\texttt{3V3}, \texttt{5V0} & PWR\\
\bottomrule
\end{tabular}
\end{center}

\subsection{Regole personalizzate (file .kicad\_dru)}
Proposta, sintassi da controllare in KiCad 10 \dv:
\begingroup\small
\begin{verbatim}
(version 1)
(rule "RF: niente via sulle reti RF a 50 ohm e RF_DIFF"
   (condition "A.Type == 'via' && (A.NetClass == 'RF_50R' || A.NetClass == 'RF_DIFF')")
   (constraint disallow via))
(rule "RF: clearance verso altre reti"
   (condition "A.NetClass == 'RF_50R' && B.Net != A.Net && B.NetName != 'GND'")
   (constraint clearance (min 0.5mm)))
(rule "Zona RF: In1 solo GND"
   (layer "In1.Cu")
   (condition "A.insideArea('ZONA_RF') && A.NetName != 'GND'")
   (constraint disallow track via zone))
\end{verbatim}
\endgroup
La terza regola richiede una \emph{rule area} chiamata \texttt{ZONA\_RF} sopra la parte RF; la seconda
non vale per la pista \net{EN_TX_ANT_SW} vicino a CTRL di IC202, che va controllata a vista.

\subsection{Impostazioni delle zone}
\begin{itemize}
\item Zona GND di F.Cu nella parte RF: clearance \SI{0.20}{mm} (uguale al gap CPWG), connessione dei pad
      \textbf{piena} (non a razze) per i pad GND dei componenti RF, larghezza minima \SI{0.2}{mm},
      rimozione delle isole.
\item Priorità della zona GND di F.Cu nella parte RF più alta delle eventuali zone 3V3 sullo stesso strato.
\end{itemize}

% =====================================================================
\section{Checklist prima dei Gerber}
\label{sec:checklist}
\subsection*{Progetto e stackup}
\begin{checklist}
\item Netclass RF\_50R, RF\_DIFF, RF\_PA\_OUT, RF\_CLK, PA\_BIAS, PWR create; tutte le reti della tabella
      \S\ref{sec:netclass} assegnate (controllo nel pannello Net Inspector).
\item Stackup in KiCad = stackup ordinato a JLCPCB; In1.Cu GND; larghezze verificate con il calcolatore
      JLCPCB (con solder mask).
\item Footprint di J201 compatibile con spessore scheda e codice JLC.
\item Via nei pad esposti di IC202, Q201, IC204 aggiunte; via di IC2 compatibili con le regole di scheda.
\end{checklist}
\subsection*{Piazzamento}
\begin{checklist}
\item IC202 con RFC verso J201; TX da RF2, RX da RF1; IC204 ruotato con il lato RF verso l'esterno
      dell'anello; nessun incrocio TX/RX.
\item J201--Q201 $\le$\,\SI{15}{mm} di pista; componenti RX in linea.
\item Rete AN643 sui pin 2/3; \net{RF_RX_P}/\net{RF_RX_N} $<$\,\SI{3}{mm}, differenza $\le$\,\SI{0.5}{mm};
      L206 tra i pad; C220 con via.
\item Rete AN648 sul pin~4, lontana dalla rete AN643, induttanze non allineate.
\item Y201 vicino a XIN, $\ge$\,\SI{15}{mm} da PA e L242; C203/C202 sul pin~4 di Y201.
\item Blocco PA copiato dalla EVB (ordine, distanze, lunghezze, larghezze); \SI{1}{mm} libero intorno a C260 e C262.
\item FL201 in linea, ingresso e uscita in direzioni opposte.
\item PA $\ge$\,\SI{15}{mm} dalla catena RX; ingresso e uscita del PA su lati opposti.
\end{checklist}
\subsection*{Sbroglio e masse}
\begin{checklist}
\item Nessuna via sulle reti RF\_50R, RF\_DIFF, \net{RF_PA_OUT}, \net{RF_PA_MATCH}; nessun cambio di strato.
\item Nessuno stub; condensatori in parallelo posati sulla pista.
\item Ogni condensatore in parallelo con la sua via; C260 e C261 con almeno 2 via; pad GND RF a connessione piena.
\item Pin~4 del PA con pista di circa \SI{1}{mm} verso la via (DS p.~7); pin~2 diretto.
\item Pin~8 (2F0) e pin~13 del PA senza pista.
\item Fence a passo $\le$\,\SI{2}{mm} su tutte le linee, $\le$\,\SI{1.5}{mm} intorno a PA e FL201; fila di
      via tra TX e RX; via su tutti i 6 pad GND di FL201.
\item In1.Cu senza piste né tagli nella zona RF; In2/B.Cu solo GND sotto le linee RF.
\item Corridoio digitale unico, a \ang{90}, sotto il tratto TX a basso livello; nessuna pista digitale a
      meno di \SI{2}{mm} dalle linee RF su F.Cu.
\item Ramo 3V3$\rightarrow$L242 dedicato, $\ge$\,\SI{1}{mm} (o pour) su strato esterno, $\ge$\,4 via per
      cambio strato; C258/C257/C256 sul lato 3V3 di L242.
\item Pad termico di IC2 con 9 via + almeno 12 via intorno + rame su B.Cu.
\end{checklist}
\subsection*{Controlli finali}
\begin{checklist}
\item DRC senza errori con le regole RF; controllo a vista delle reti RF in 3D e sul layer F.Cu da solo.
\item Confronto netlist--PCB (\emph{Update PCB from Schematic} senza differenze; R243 = \SI{220}{\ohm}).
\item Piazzole di misura (GND senza mask) accanto a C230, C255, C212, C213 se volute.
\item Serigrafia: riferimenti leggibili di C260, C262, C219, C220, L205, L206, C702, C704; nessuna
      serigrafia sui pad.
\item Revisione del layout RF con questa checklist firmata da un secondo progettista.
\end{checklist}

% =====================================================================
\appendix
\section{Componenti RF: footprint e note di piazzamento}
\label{app:componenti}
{\footnotesize
\renewcommand{\arraystretch}{1.08}
\begin{xltabular}{\linewidth}{@{}L{11mm}L{23mm}L{35mm}Y@{}}
\toprule
\textbf{Ref} & \textbf{Valore} & \textbf{Footprint} & \textbf{Rete / note di piazzamento}\\
\midrule\endhead
\midrule\multicolumn{4}{r}{\footnotesize\itshape continua\ldots}\\\endfoot
\bottomrule\endlastfoot
\multicolumn{4}{@{}l}{\textbf{Antenna e switch}}\\
J201 & MMCX & MMCX Molex 73415-0961 0.8mm-PCB & bordo scheda; 3 via per pad GND; vedi P05\\
C206 & 100p & C\_0603 & \net{RF_ANT}--\net{RF_SW_COM}, subito dopo J201\\
IC202 & HMC349AMS8GE & SOP65P490X110-9N & RFC verso J201; via nel pad esposto (da aggiungere)\\
C211 & 100n & C\_0402 & VDD pin~1 (5V0), via sul pad\\
C207 & 100p & C\_0603 & sul pin RF1 (RX)\\
C208 & 100p & C\_0603 & sul pin RF2 (TX)\\
\multicolumn{4}{@{}l}{\textbf{RX prima del LNA}}\\
C210 & 6.8p & C\_0805 & \net{RF_RX_F1} verso GND, via propria\\
L201 & 22n & L\_0805 & serie, asse lungo la pista\\
C209 & 6.8p & C\_0805 & \net{RF_RX_F2} verso GND, via propria\\
C212 & 100p & C\_0603 & DC block sul pin~2 di Q201; piazzola GND di misura vicina\\
Q201 & TQP3M9036 & SON50P200X200X95-9N-D & via nel pad esposto (da aggiungere); NC a GND\\
R202 & 0R & R\_0402 & pin~6, lato opposto all'ingresso\\
R203 & DNP & R\_0402 & pin~6 verso GND\\
L202 & 120n & L\_0402 HandSolder & choke uscita LNA, perpendicolare a L203\\
C214 & 100p & C\_0402 & lato 3V3 di L202, via sul pad\\
C215 & 1u & C\_0402 & dopo C214\\
\multicolumn{4}{@{}l}{\textbf{RX dopo il LNA e AN643}}\\
C213 & 100p & C\_0603 & serie, uscita LNA\\
C217 & 5.6p & C\_0603 & \net{RF_RX_F3} verso GND\\
L203 & 22n & L\_0805 & serie, non affiancata a L204\\
C216 & 8.2p & C\_0805 & \net{RF_RX_F4} verso GND\\
L204 & 22n & L\_0805 & serie\\
C218 & 5.6p & C\_0603 & \net{RF_RX_MATCH} verso GND\\
C219 & 2.2p & C\_0603 & serie verso RXP; simmetrico a L205; spazio taratura\\
L205 & 56n & L\_0603 & serie verso RXN\\
L206 & 62n & L\_0603 & tra RXP e RXN, a meno di 1\,mm dai pin\\
C220 & 4.7p & C\_0805 & RXN verso GND, via sul pad\\
\multicolumn{4}{@{}l}{\textbf{Si4464 e TCXO}}\\
IC204 & SI4464-B1B-FMR & QFN50P400X400X90-21N-D & ruotato \ang{180}; via nel pad esposto (da aggiungere)\\
C201 & 100n & C\_0402 & VDD pin~6/8, via sul pad\\
Y201 & ECS-TXO-20CSMV4-320 & ECSTXO20CSMV4320AYTR & vicino a XIN, $\ge$\,15\,mm dal PA\\
C202 & 1u & C\_0402 & VDD di Y201\\
C203 & 100n & C\_0402 & VDD di Y201, più vicino al pin\\
C204 & 10n & C\_0402 & accoppiamento AC, in linea\\
R201 & 33 & R\_0402 & in serie verso XIN\\
TP701 & -- & TestPoint D1.0mm & sulla pista \net{RF_XIN}, senza stub\\
\multicolumn{4}{@{}l}{\textbf{TX: AN648 e attenuatore}}\\
L703 & 220n & L\_0603 & choke sul pin~4; 3V3 con via vicino a C201\\
C702 & 9.1p & C\_0402 & serie sul pin~4\\
L702 & 33n & L\_0805 & serie, asse perpendicolare a L205/L206\\
C704 & 15p & C\_0402 & \net{RF_TX_M2} verso GND\\
C701 & 100p & C\_0603 & DC block verso il pad\\
R245 & 100 & R\_0603 & shunt ingresso pad, \SI{49}{mW} max\\
R246 & 68 & R\_0603 & serie\\
R247 & 100 & R\_0402 & shunt uscita pad\\
R248 & 180 & R\_0402 & shunt uscita pad (R1 EVB)\\
C255 & 56p & C\_0603 & DC block ingresso PA (C1 EVB), \SI{0.9}{mm} dal pin~3\\
\multicolumn{4}{@{}l}{\textbf{PA QPA9510}}\\
IC2 & QPA9510TR7 & QFN50P300X300X90-17N-D & 9 via nel pad (foro 0,25); blocco EVB ruotato \ang{180}\\
R240 & 0 & R\_0402 HS & 3V3$\rightarrow$\net{PA_VCC1_F} (L7 EVB, lì DNP)\\
L240 & 12n & L\_0603 HS & \net{PA_VCC1}, posizione EVB (L1)\\
C244 & 330p & C\_0402 HS & \net{PA_VCC1_F} (C19 EVB)\\
C240, C241 & 10n, 1n & C\_0402 HS & 3V3 lato VCC1 (C2, C3 EVB)\\
R241 & 0 & R\_0603 HS & 3V3$\rightarrow$\net{PA_VCC2_F} (L2 EVB)\\
L241 & 6.8n & L\_0402 HS & \net{PA_VCC2}, posizione EVB (L6)\\
C248 & 330p & C\_0402 HS & \net{PA_VCC2_F} (C20 EVB)\\
C245, C246, C247 & 10u, 10n, 1n & C\_0805 HS / C\_0402 HS & 3V3 lato VCC2 (C21, C5, C6 EVB)\\
C251 & 100p & C\_0402 HS & VCC pin~14 (C15 EVB)\\
R242, R243 & 33, 220 & R\_0402 HS & partitore APC, vicino ai pin 15/16\\
C252 & 10n & C\_0402 HS & APC verso GND (C16 EVB)\\
Q701 & FDN308P & SOT-23 & interruttore APC, dentro l'anello\\
L242 & 1u NRS2012 & L\_Taiyo-Yuden\_NR-20xx & choke stadio finale; pista 3V3 $\ge$\,1\,mm\\
C256, C257, C258 & 10u, 1n, 47p & C\_0805 HS / C\_0402 HS & lato 3V3 di L242, 47p più vicino\\
C260 & 27p & C\_0603 & nodo \net{RF_PA_OUT}, $\ge$\,2 via, 1\,mm libero\\
L243 & 2.7n & L\_0402 & serie, posizione EVB (L4)\\
C261 & 22p & C\_0402 & nodo \net{RF_PA_MATCH}, $\ge$\,2 via\\
C262 & 2p & C\_0402 & taratura (C25 EVB), 1\,mm libero\\
C230 & 56p & C\_0603 & DC block verso FL201 (C12 EVB)\\
\multicolumn{4}{@{}l}{\textbf{Filtro}}\\
FL201 & LFCG-490+ & IC:LFCG490 & in linea; via su 6 pad GND; footprint da confrontare con PL-429\\
\end{xltabular}}
\noindent{\footnotesize HS = footprint \enquote{HandSolder}. Valori e footprint dalla netlist rev.~1.45.}

% =====================================================================
\section{Punti da verificare e incongruenze trovate}
\label{app:verifiche}

\subsection{Da verificare}
\begin{enumerate}
\item Larghezze delle linee con il calcolatore JLCPCB e con la solder mask (effetto stimato 1--3\,\si{\ohm}).
\item Lunghezze e larghezze delle linee della EVB QPA9510 (lette da una figura a bassa risoluzione, unità
      della larghezza \enquote{19,8} non chiara) e substrato della EVB: servono i Gerber Qorvo.
\item Effetto della larghezza EVB (\SI{0.5}{mm}) sul nostro stackup nella rete del PA (circa \SI{26}{\ohm}
      invece di \SI{50}{\ohm}): taratura di C260/C262 al banco (M-06).
\item Footprint di J201 (\SI{0.8}{mm}) contro spessore della scheda e codice DreamLNK (P05).
\item Footprint \texttt{IC:LFCG490} contro il layout consigliato Mini-Circuits PL-429.
\item Numero e diametro delle via nei pad esposti di IC204 e IC202 (AN629 e disegni non consultati); pasta a finestre.
\item Disponibilità di via riempite e chiuse (via-in-pad) sul 4 strati JLC.
\item Rame nel foro delle via (\SI{20}{\micro m} assunto) e conto termico del PA con l'area di rame reale.
\item Regole attribuite ad AN629 (orientamento delle induttanze, disaccoppiamento del Si4464): il PDF
      Silicon Labs non era accessibile.
\item Distanze indicative senza fonte diretta: TCXO--XIN $\le$\,\SI{5}{mm}; L242/L243 da L702/L703
      $\ge$\,\SI{5}{mm}; buck $\ge$\,\SI{20}{mm} dalla zona RF; area di rame B.Cu del PA.
\item Sintassi delle regole \texttt{.kicad\_dru} in KiCad 10.
\end{enumerate}

\subsection{Incongruenze (segnalate, non corrette)}
\begin{enumerate}
\item \textbf{Netclass assenti}: \texttt{SYMPHONY.kicad\_pro} ha solo \texttt{Default}, nessuna assegnazione
      e nessun pattern; il log (1.12-02, 1.13-02) le dà come inserite. Nessun backup in
      \texttt{report\_src} le contiene.
\item \textbf{Stackup del file PCB}: \texttt{SYMPHONY.kicad\_pcb} ha 6 strati con In1 di segnale sotto F.Cu;
      il log \S24.1 parla di 4 strati con In1 GND.
\item \textbf{Spessori del \S24.1}: core \SI{0.535}{mm} con prepreg \SI{0.1}{mm} non dà \SI{1.6}{mm}; il
      JLC04161H-3313 ha core \SI{1.265}{mm}.
\item \textbf{Schema di posizionamento \S24.3}: con l'ordine dei pin di IC202 (RF1 = RX in basso, RF2 = TX in
      alto) e di IC204 (RXP/RXN sopra TX), la disposizione del disegno porta a un incrocio TX/RX. Soluzioni:
      IC204 ruotato con i pin non RF dentro l'anello (questa guida), oppure scambio RF1/RF2 con CTRL preso
      dal segnale invertito (modifica di schema).
\item \textbf{L-01 punti 1--2}: copiare le larghezze EVB su uno stackup diverso non copia l'impedenza delle
      linee (EVB circa \SI{50}{\ohm}, nostro 3313 circa \SI{26}{\ohm} a \SI{0.5}{mm}).
\item \textbf{Netclass RF\_PA\_OUT a \SI{1.0}{mm}} (\S24.2) contro \enquote{stesse larghezze della EVB} (L-01):
      a \SI{1.0}{mm} il tratto vale circa \SI{15}{\ohm}.
\item \textbf{IPC interno nel \S24.4}: \SI{1.6}{mm} per \SI{1.7}{A} vale con rame interno da 1 oz; con il
      rame interno JLC da 0,5 oz servono circa \SI{3.7}{mm}.
\item \textbf{Footprint senza via termiche}: pad esposti di IC204 (Si4464), Q201 (TQP3M9036, il datasheet
      le richiede) e IC202 (HMC349A) senza via.
\item \textbf{Via di IC2 sotto il minimo di scheda}: foro \SI{0.25}{mm} contro \texttt{min\_through\_hole\_diameter}
      \SI{0.3}{mm}.
\item \textbf{L703 senza disaccoppiamento proprio} sul lato 3V3 (solo C201 sul VDD del Si4464).
\item \textbf{Frequenza del TCXO}: il disegno del \S24.3 scrive \enquote{TCXO 30 MHz}; Y201 è da
      \SI{32}{MHz} (codice \texttt{-320}).
\item \textbf{Simulazione Qorvo}: nel modello di p.~12 L4 sembra \SI{2.4}{nH} (LQG15HS), nella BOM EVB è
      \SI{2.7}{nH} (lettura della figura \dv).
\item \textbf{Log 001 \S5}: \enquote{EN\_TX\_ANT\_SW = 0 vuol dire switch su RF1 (TX)}, mentre nella netlist
      RF1 è la porta RX (\net{RF_SW_RX}). Riguarda firmware e logica, non il layout; da ricontrollare con la
      tabella di verità HMC349A (DS p.~8).
\end{enumerate}

% =====================================================================
\section{Fonti}
\label{app:fonti}
{\small
\begin{enumerate}[label={[\arabic*]}]
\item Netlist KiCad di \texttt{SYMPHONY.kicad\_sch}, rev.~1.45, esportata con \texttt{kicad-cli} 10.0.8 il
      04/11/2026 (\texttt{docs/layout\_rf/net.xml}).
\item \texttt{SYMPHONY.kicad\_pro} e \texttt{SYMPHONY.kicad\_pcb} (stato al 04/11/2026), footprint in
      \texttt{UHF\_LIB/FOOTPRINT/IC.pretty} e libreria KiCad 10 \texttt{Connector\_Coaxial}.
\item Log hardware SYMPHONY V2, \texttt{report/entries}: 001 (\S3, \S4, \S8--10), 007 (\S21), 009
      (\S23), 010 (\S24), 011 (\S25), 023 (\S37), 041, 043 (L-01, M-06).
\item Revisione indipendente del 04/11/2026, \texttt{DESIGN/SYMPHONY\_V2\_review\_finale\_2026-11-04.pdf}:
      P05, P08, P10, P11, \S2.2.
\item Qorvo, \emph{QPA9510 3V General Purpose/GSM Power Amplifier}, data sheet 23/06/2025: p.~6 (EVB 900 MHz),
      p.~7 (descrizione dei pad), p.~8 (PCB mounting pattern).
\item Qorvo, \emph{QPA9510 410--470 MHz EVB Test Summary}, 30/06/2025: p.~2 (foto e schema EVB), p.~3 (BOM),
      p.~6--7 (guadagno e P1dB), p.~11--12 (modello e impedenze simulate).
\item Analog Devices, \emph{HMC349AMS8G}, data sheet rev.~F: p.~5 (pin), p.~8 (teoria di funzionamento),
      p.~9 (scheda di valutazione).
\item Qorvo, \emph{TQP3M9036}, data sheet 27/11/2018: p.~10 (pin e stackup EVB), p.~11 (PCB mounting pattern).
\item Silicon Labs, \emph{Si4464/63/61/60}, data sheet rev.~1.4: p.~38 (TCXO su XIN), p.~46--47 (pin, pad esposto).
\item Mini-Circuits, \emph{LFCG-490+}, data sheet: p.~1--2 (prestazioni, potenza), p.~3 (pad, layout PL-429).
\item ECS, \emph{ECS-TXO-20CSMV4}, data sheet: pin, livello d'uscita, codici di ordinazione.
\item JLCPCB, \emph{Impedance / stackup} (\texttt{jlcpcb.com/impedance}) e \emph{PCB capabilities}
      (\texttt{jlcpcb.com/capabilities/pcb-capabilities}), consultati il 04/11/2026.
\item Silicon Labs AN629, AN643, AN648: non accessibili durante la stesura (HTTP 403); valori AN643/AN648 dal log
      [3]; una sintesi di AN629 è in \emph{Microwaves \& RF}, \enquote{Investigating ICs for Tomorrow's Radios}.
\item IPC-2221, formula della corrente nei conduttori ($I=k\Delta T^{0{,}44}A^{0{,}725}$).
\item Formule di linea: Hammerstad e Jensen (microstriscia, 1980); trasformazione conforme per CPWG con
      approssimazione di Hilberg; induttanza della via $L\approx0{,}2h[\ln(4h/d)+1]$\,nH.
\end{enumerate}}

\end{document}
